% Included by main.tex immediately after Appendix E.
\clearpage
\begingroup
\edef\suppSavedSection{\number\value{section}}
\edef\suppSavedEquation{\number\value{equation}}
\edef\suppSavedTable{\number\value{table}}
\newgeometry{margin=25mm,top=23mm,bottom=23mm}
\fontsize{10}{12}\selectfont
\let\small\footnotesize
\renewcommand{\thesection}{S\arabic{section}}
\renewcommand{\theHsection}{supp.\arabic{section}}
\renewcommand{\thetable}{S\arabic{table}}
\renewcommand{\theHtable}{supp.\arabic{table}}
\setcounter{section}{0}
\setcounter{equation}{0}
\setcounter{table}{0}
\makeatletter
\@fleqntrue
\@mathmargin=0pt
\makeatother
\let\suppaligned\aligned
\let\endsuppaligned\endaligned
\section*{Supplementary calculations}
\label{sec:supplementary-calculations}
\addcontentsline{toc}{section}{Supplementary calculations}

\subsection*{Relation to the proof}
\noindent\begin{tabular}{@{}>{\raggedright\arraybackslash}p{0.51\linewidth}
                       >{\raggedright\arraybackslash}p{0.43\linewidth}@{}}
\toprule Supplementary outputs & Results in the proof\\\midrule
Section~\ref{supp:frame}: derivative matrices, compatibility polynomials
and reconstruction & Proposition~\ref{fnd:finiteclosure};
Proposition~\ref{fnd:reconstruction}\\[7pt]
Section~\ref{supp:normalization}: ratio equations, normalization factors
and similarity orbits & Lemmas~\ref{fnd:constantratio}
and~\ref{fnd:muwall}; Theorem~\ref{fnd:lowrankexclusion}\\[7pt]
Section~\ref{supp:isometries}: connection fibres and cylindrical
linearization & Theorem~\ref{iso:theorem};
Lemma~\ref{iso:compactness}; Lemma~\ref{iso:rr-linearization}\\[7pt]
Section~\ref{supp:dilations}: homothety equations, principal-direction
normal equations and carrier matrices & Theorem~\ref{dil:whole};
Lemma~\ref{dil:cyl-bound}; Lemma~\ref{dil:carrier}\\[7pt]
Section~\ref{supp:boundary}: local ideals, normal determinants
and ordinary boundary obstructions & Lemma~\ref{r3:boundaryreduction};
Propositions~\ref{r3:ordinaryexclusion}
and~\ref{r3:dominantfinite}\\[7pt]
Section~\ref{supp:curves}: characteristic pencils, direction factors
and amplitude conics & Proposition~\ref{r3:curvefibre}\\[7pt]
Sections~\ref{supp:elimination} and~\ref{supp:momentcoefficients}:
moment matrix, quadratic norm and slope resultants &
Propositions~\ref{r3:directioneliminants}
and~\ref{r3:curvefibre}\\[7pt]
Section~\ref{supp:classical}: deformation gradients and pressures &
Proposition~\ref{completion:sufficiency}\\
\bottomrule
\end{tabular}

\section{The finite frame system}\label{supp:frame}

We use the notation and derivative ordering of
Section~\ref{sec:geometry}. In particular $h_i=p_{ii}$,
$\boldsymbol\tau=(\tau_0,\ldots,\tau_5)^{\top}$ and
$\boldsymbol u=(u_0,u_1)^{\top}$. Coefficients lie in
$\mathbb Q(a,b)$, where
\[
 a,b>0,\qquad ab(a-1)(b-1)(a-b)\ne0.
\]
The raw state has coordinate weights $(1^9,2^6,3^2)$.
All row and column indices below start at one.

\subsection{The three derivative matrices}

The curvature and pressure expressions define $M^{(0)}$; commutators
of the frame coefficients define $M^{(1)}$; commutators of the
retained first derivatives define $M^{(2)}$. Their dimensions and
unknown columns are
\[
\begin{array}{ccl}
 M^{(0)}\in\operatorname{Mat}_{24\times27}(\mathbb Q(a,b)),
 & & (D_ip_{jk})_{i,j,k},\\
 M^{(1)}\in\operatorname{Mat}_{27\times18}(\mathbb Q(a,b)),
 & & (D_i\tau_r)_{i=1,2,3;\,r=0,\ldots,5},\\
 M^{(2)}\in\operatorname{Mat}_{18\times6}(\mathbb Q(a,b)),
 & & (D_iu_r)_{i=1,2,3;\,r=0,1}.
\end{array}
\]
The indices $j,k$ are lexicographic. The commutator pairs are
$(2,3),(3,1),(1,2)$, followed by the coordinate index. The exact
coefficient definitions are those of Appendix~\ref{appf:finiteconstruction}.
Tables~\ref{supp:matrix-M0}--\ref{supp:matrix-M2} give the resulting matrices.

\begingroup\small
\setlength{\tabcolsep}{3pt}
\renewcommand{\arraystretch}{1.16}
\begin{longtable}{@{}r>{\raggedright\arraybackslash}p{0.91\linewidth}@{}}
\caption{The first derivative matrix $M^{(0)}$. An entry $(j,c)$ in row $i$ means that the $(i,j)$ coefficient is $c$. Unlisted entries are zero.}\label{supp:matrix-M0}\\
\toprule Row & Nonzero column entries\\\midrule\endfirsthead
\toprule Row & Nonzero column entries\\\midrule\endhead
\bottomrule\endfoot
1 & \((16,-1)\)\allowbreak\quad \((22,1)\)\\
2 & \((7,1)\)\allowbreak\quad \((19,-1 \cdot 2^{-1})\)\allowbreak\quad \((23,2^{-1})\)\allowbreak\quad \((27,2^{-1})\)\\
3 & \((4,-1)\)\allowbreak\quad \((10,2^{-1})\)\allowbreak\quad \((14,-1 \cdot 2^{-1})\)\allowbreak\quad \((18,-1 \cdot 2^{-1})\)\\
4 & \((17,-1)\)\allowbreak\quad \((19,-1 \cdot 2^{-1})\)\allowbreak\quad \((23,2^{-1})\)\allowbreak\quad \((27,-1 \cdot 2^{-1})\)\\
5 & \((8,1)\)\allowbreak\quad \((20,-1)\)\\
6 & \((1,2^{-1})\)\allowbreak\quad \((5,-1 \cdot 2^{-1})\)\allowbreak\quad \((9,2^{-1})\)\allowbreak\quad \((11,1)\)\\
7 & \((10,2^{-1})\)\allowbreak\quad \((14,2^{-1})\)\allowbreak\quad \((18,-1 \cdot 2^{-1})\)\allowbreak\quad \((24,1)\)\\
8 & \((1,-1 \cdot 2^{-1})\)\allowbreak\quad \((5,-1 \cdot 2^{-1})\)\allowbreak\quad \((9,2^{-1})\)\allowbreak\quad \((21,-1)\)\\
9 & \((6,-1)\)\allowbreak\quad \((12,1)\)\\
10 & \((16,\left(-1\right) \cdot b^{-1})\)\allowbreak\quad \((22,a^{-1})\)\\
11 & \((7,b^{-1})\)\allowbreak\quad \((19,-1 \cdot 2^{-1})\)\allowbreak\quad \((23,\left(\frac{1}{2}\right) \cdot a^{-1})\)\allowbreak\quad \((27,\left(\frac{1}{2}\right) \cdot b^{-1})\)\\
12 & \((4,\left(-1\right) \cdot a^{-1})\)\allowbreak\quad \((10,2^{-1})\)\allowbreak\quad \((14,\left(-\frac{1}{2}\right) \cdot a^{-1})\)\allowbreak\quad \((18,\left(-\frac{1}{2}\right) \cdot b^{-1})\)\\
13 & \((17,\left(-1\right) \cdot b^{-1})\)\allowbreak\quad \((19,-1 \cdot 2^{-1})\)\allowbreak\quad \((23,\left(\frac{1}{2}\right) \cdot a^{-1})\)\allowbreak\quad \((27,\left(-\frac{1}{2}\right) \cdot b^{-1})\)\\
14 & \((8,b^{-1})\)\allowbreak\quad \((20,-1)\)\\
15 & \((1,2^{-1})\)\allowbreak\quad \((5,\left(-\frac{1}{2}\right) \cdot a^{-1})\)\allowbreak\quad \((9,\left(\frac{1}{2}\right) \cdot b^{-1})\)\allowbreak\quad \((11,1)\)\\
16 & \((10,2^{-1})\)\allowbreak\quad \((14,\left(\frac{1}{2}\right) \cdot a^{-1})\)\allowbreak\quad \((18,\left(-\frac{1}{2}\right) \cdot b^{-1})\)\allowbreak\quad \((24,a^{-1})\)\\
17 & \((1,-1 \cdot 2^{-1})\)\allowbreak\quad \((5,\left(-\frac{1}{2}\right) \cdot a^{-1})\)\allowbreak\quad \((9,\left(\frac{1}{2}\right) \cdot b^{-1})\)\allowbreak\quad \((21,-1)\)\\
18 & \((6,\left(-1\right) \cdot a^{-1})\)\allowbreak\quad \((12,1)\)\\
19 & \((11,(b - 1))\)\allowbreak\quad \((13,(a - b))\)\allowbreak\quad \((21,(a - 1))\)\allowbreak\quad \((25,\left(-1\right) \cdot (a - b))\)\\
20 & \((2,\left(-1\right) \cdot (b - 1))\)\allowbreak\quad \((4,\left(-1\right) \cdot (a - b))\)\allowbreak\quad \((24,\left(-1\right) \cdot (a - 1))\)\allowbreak\quad \((26,(b - 1))\)\\
21 & \((3,\left(-1\right) \cdot (a - 1))\)\allowbreak\quad \((7,(a - b))\)\allowbreak\quad \((15,(a - 1))\)\allowbreak\quad \((17,\left(-1\right) \cdot (b - 1))\)\\
22 & \((11,\left(-1\right) \cdot b^{-1} \cdot (b - 1))\)\allowbreak\quad \((13,\left(-1\right) \cdot b^{-1} \cdot a^{-1} \cdot (a - b))\)\allowbreak\quad \((21,\left(-1\right) \cdot a^{-1} \cdot (a - 1))\)\allowbreak\quad \((25,b^{-1} \cdot a^{-1} \cdot (a - b))\)\\
23 & \((2,b^{-1} \cdot (b - 1))\)\allowbreak\quad \((4,b^{-1} \cdot a^{-1} \cdot (a - b))\)\allowbreak\quad \((24,a^{-1} \cdot (a - 1))\)\allowbreak\quad \((26,\left(-1\right) \cdot b^{-1} \cdot (b - 1))\)\\
24 & \((3,a^{-1} \cdot (a - 1))\)\allowbreak\quad \((7,\left(-1\right) \cdot b^{-1} \cdot a^{-1} \cdot (a - b))\)\allowbreak\quad \((15,\left(-1\right) \cdot a^{-1} \cdot (a - 1))\)\allowbreak\quad \((17,b^{-1} \cdot (b - 1))\)\\
\end{longtable}
\endgroup

\begingroup\small
\setlength{\tabcolsep}{3pt}
\renewcommand{\arraystretch}{1.16}
\begin{longtable}{@{}r>{\raggedright\arraybackslash}p{0.91\linewidth}@{}}
\caption{The first prolongation matrix $M^{(1)}$. An entry $(j,c)$ in row $i$ means that the $(i,j)$ coefficient is $c$. Unlisted entries are zero.}\label{supp:matrix-M1}\\
\toprule Row & Nonzero column entries\\\midrule\endfirsthead
\toprule Row & Nonzero column entries\\\midrule\endhead
\bottomrule\endfoot
1 & \((9,\left(-1\right) \cdot (b - 1)^{-1} \cdot a^{-1} \cdot (a - b))\)\allowbreak\quad \((16,\left(-2\right) \cdot (b - 1)^{-1} \cdot a^{-1} \cdot (a - b))\)\\
2 & \((14,1)\)\\
3 & \((8,1)\)\\
4 & \((14,1)\)\allowbreak\quad \((17,-1)\)\\
5 & \((9,1)\)\\
6 & \((10,1)\)\allowbreak\quad \((13,-1)\)\\
7 & \((11,1)\)\\
8 & \((12,1)\)\allowbreak\quad \((15,\left(-\frac{1}{2}\right) \cdot (b - 1)^{-1} \cdot a^{-1} \cdot b \cdot (a - 1))\)\\
9 & \((16,\left(-2\right) \cdot (b - 1)^{-1} \cdot a^{-1} \cdot b \cdot (a - 1))\)\\
10 & \((3,(b - 1)^{-1} \cdot a^{-1} \cdot (a - b))\)\\
11 & \((16,-1)\)\allowbreak\quad \((18,1)\)\\
12 & \((2,-1)\)\allowbreak\quad \((13,1)\)\allowbreak\quad \((15,\left(-\frac{1}{2}\right) \cdot (b - 1)^{-1} \cdot a^{-1} \cdot b \cdot (a - 1))\)\\
13 & \((16,-1)\)\\
14 & \((3,-1)\)\allowbreak\quad \((14,\left(2\right) \cdot (a - b)^{-1} \cdot (b - 1) \cdot a)\)\\
15 & \((4,-1)\)\\
16 & \((5,-1)\)\allowbreak\quad \((15,\left(-\frac{1}{2}\right) \cdot (b - 1)^{-1} \cdot a^{-1} \cdot b \cdot (a - 1))\)\\
17 & \((6,-1)\)\\
18 & \((14,\left(2\right) \cdot (a - b)^{-1} \cdot b \cdot (a - 1))\)\\
19 & \((4,\left(2\right) \cdot (b - 1)^{-1} \cdot a^{-1} \cdot (a - b))\)\\
20 & \((2,-1)\)\allowbreak\quad \((10,1)\)\allowbreak\quad \((12,-1)\)\\
21 & \((7,-1)\)\allowbreak\quad \((9,\left(\frac{1}{2}\right) \cdot (b - 1)^{-1} \cdot a^{-1} \cdot b \cdot (a - 1))\)\\
22 & \((2,-1)\)\allowbreak\quad \((5,1)\)\allowbreak\quad \((10,1)\)\\
23 & \((8,\left(-2\right) \cdot (a - b)^{-1} \cdot (b - 1) \cdot a)\)\\
24 & \((1,1)\)\\
25 & \((9,\left(\frac{1}{2}\right) \cdot (b - 1)^{-1} \cdot a^{-1} \cdot b \cdot (a - 1))\)\\
26 & \((3,\left(\frac{1}{2}\right) \cdot (b - 1)^{-1} \cdot a^{-1} \cdot b \cdot (a - 1))\)\\
27 & \((4,\left(2\right) \cdot (b - 1)^{-1} \cdot a^{-1} \cdot b \cdot (a - 1))\)\allowbreak\quad \((8,\left(-2\right) \cdot (a - b)^{-1} \cdot b \cdot (a - 1))\)\\
\end{longtable}
\endgroup

\begingroup\small
\setlength{\tabcolsep}{3pt}
\renewcommand{\arraystretch}{1.16}
\begin{longtable}{@{}r>{\raggedright\arraybackslash}p{0.91\linewidth}@{}}
\caption{The second prolongation matrix $M^{(2)}$. An entry $(j,c)$ in row $i$ means that the $(i,j)$ coefficient is $c$. Unlisted entries are zero.}\label{supp:matrix-M2}\\
\toprule Row & Nonzero column entries\\\midrule\endfirsthead
\toprule Row & Nonzero column entries\\\midrule\endhead
\bottomrule\endfoot
1 & \((3,1)\)\\
3 & \((4,1)\)\\
4 & \((5,-1)\)\\
6 & \((6,\left(-\frac{1}{2}\right) \cdot (b - 1)^{-1} \cdot a^{-1} \cdot b \cdot (a - 1))\)\\
7 & \((1,-1)\)\\
8 & \((5,1)\)\allowbreak\quad \((6,\left(-\frac{1}{2}\right) \cdot (b - 1)^{-1} \cdot a^{-1} \cdot b \cdot (a - 1))\)\\
9 & \((2,-1)\)\\
11 & \((6,\left(-\frac{1}{2}\right) \cdot (b - 1)^{-1} \cdot a^{-1} \cdot b \cdot (a - 1))\)\\
14 & \((3,-1)\)\allowbreak\quad \((4,\left(\frac{1}{2}\right) \cdot (b - 1)^{-1} \cdot a^{-1} \cdot b \cdot (a - 1))\)\\
16 & \((1,1)\)\\
17 & \((4,\left(\frac{1}{2}\right) \cdot (b - 1)^{-1} \cdot a^{-1} \cdot b \cdot (a - 1))\)\\
18 & \((2,\left(\frac{1}{2}\right) \cdot (b - 1)^{-1} \cdot a^{-1} \cdot b \cdot (a - 1))\)\\
\end{longtable}
\endgroup

For the row and column selections in Appendix~\ref{appf:finiteconstruction},
the selected determinants are

\begin{equation}\label{supp:fnd:firstminor}
\det M^{(0)}_{I_0,J_0}=-\frac{(a-1)^4(a-b)^5(b-1)^6}{32a^6b^8},
\qquad M^{(0)}Q-r^{(0)}=0.
\end{equation}

\begin{equation}\label{supp:fnd:secondminor}
\det M^{(1)}_{I_1,J_1}=-\frac{2(a-b)^2}{a^2(b-1)^2}.
\end{equation}

\begin{equation}\label{supp:appf:minor2}
\det M^{(2)}_{I_2,\{1,\ldots,6\}}=\frac{b(a-1)}{2a(b-1)}.
\end{equation}

They determine the three raw state fields $\mathcal Q_i$. Their
weights satisfy

\begin{equation}\label{supp:appf:weightidentity}
 \sum_{\ell=1}^{17}d_\ell \mathsf S_\ell
       \frac{\partial\mathcal Q_i^j}{\partial \mathsf S_\ell}
 =(d_j+1)\mathcal Q_i^j.
\end{equation}

\subsection{Cubic and quartic compatibility}

The first prolongation has three monic cubic residuals, in the
order used to define the quartic equations.

\begin{equation}\label{supp:fnd:cubic1}
\begin{aligned}
C_1={}&p_{11}^{2} p_{23} - p_{11}^{2} p_{32} + \frac{p_{21} \tau_{3} \left(a - b\right)}{a b - a} + \frac{p_{22} \tau_{1} \left(a - b\right)}{a b - a}\\
&{}+(\frac{p_{21}^{2} p_{23} \left(- a + b\right)}{a b - a} + \frac{p_{22} \tau_{4} \left(- a + b\right)}{a b - a})\\
&{}+(\frac{p_{21}^{2} p_{32} \left(- a + b\right)}{a b - a} + \frac{p_{31} \tau_{2} \left(- a + b\right)}{2 a b - 2 a})\\
&{}+(\frac{p_{11} p_{32} p_{33} \left(a - b\right)}{a b - b} + \frac{p_{33} \tau_{4} \left(a - b\right)}{a b - b})\\
&{}+(\frac{p_{21} p_{22} p_{31} \left(2 a - 2 b\right)}{a b - a} + \frac{p_{21} p_{31} p_{33} \left(- a + b\right)}{a b - b})
\end{aligned}
\end{equation}

\begin{equation}\label{supp:fnd:cubic2}
\begin{aligned}
C_2={}&p_{11} p_{12} p_{32} - p_{11} \tau_{3} + p_{11} \tau_{5} - p_{12}^{2} p_{13} - p_{12}^{2} p_{31} - p_{12} \tau_{1}\\
&{}+(\frac{p_{13} p_{22}^{2} \left(a - b\right)}{a b - a} + \frac{p_{22}^{2} p_{31} \left(- a + b\right)}{a b - a})\\
&{}+(\frac{p_{32} \tau_{2} \left(- a + b\right)}{2 a b - 2 a} + \frac{p_{33} \tau_{5} \left(- a + b\right)}{a b - b})\\
&{}+(\frac{p_{12} p_{32} p_{33} \left(- a + b\right)}{a b - b} + \frac{p_{21} p_{22} p_{32} \left(a - b\right)}{a b - a})\\
&{}+(\frac{p_{22} p_{31} p_{33} \left(a - b\right)}{a b - b})
\end{aligned}
\end{equation}

\begin{equation}\label{supp:fnd:cubic3}
\begin{aligned}
C_3={}&p_{11} p_{12} p_{33} - p_{11} p_{13} p_{32} - p_{11} p_{23} p_{31} - p_{11} p_{31} p_{32} - p_{11} \tau_{0} - p_{12} p_{13}^{2} - p_{13} \tau_{1}\\
&{}+(2 p_{11} p_{13} p_{23} + \frac{p_{11} \tau_{2} \left(a b - b\right)}{2 a b - 2 a} - p_{12} p_{13} p_{31})\\
&{}+(\frac{p_{12} p_{33}^{2} \left(- a + b\right)}{a b - b} + \frac{p_{21} p_{23}^{2} \left(a - b\right)}{a b - a})\\
&{}+(\frac{p_{21} p_{33}^{2} \left(a - b\right)}{a b - b} + \frac{p_{22} \tau_{0} \left(- a + b\right)}{a b - a})\\
&{}+(\frac{p_{11} p_{21} p_{22} \left(- a b + b\right)}{a b - a} + \frac{p_{23} \tau_{3} \left(- a + b\right)}{a b - a})\\
&{}+(\frac{p_{21} p_{22} p_{33} \left(- a + b\right)}{a b - a} + \frac{p_{21} p_{23} p_{32} \left(a - b\right)}{a b - a})
\end{aligned}
\end{equation}

The twenty-one quartics consist of the twelve nonzero residuals
of the second prolongation, followed by
$\mathcal Q_1C_1,\mathcal Q_2C_1,\mathcal Q_3C_1$, and the analogous
triples for $C_2,C_3$. Write them as
$M\boldsymbol u-\boldsymbol r$. Table~\ref{supp:matrix-Mu} gives all
coefficients of the $21\times2$ matrix $M$.

\begingroup\small
\setlength{\tabcolsep}{3pt}
\renewcommand{\arraystretch}{1.16}
\begin{longtable}{@{}r>{\raggedright\arraybackslash}p{0.91\linewidth}@{}}
\caption{The coefficient matrix of the retained second derivatives. An entry $(j,c)$ in row $i$ means that the $(i,j)$ coefficient is $c$. Unlisted entries are zero.}\label{supp:matrix-Mu}\\
\toprule Row & Nonzero column entries\\\midrule\endfirsthead
\toprule Row & Nonzero column entries\\\midrule\endhead
\bottomrule\endfoot
1 & \((1,\left(\frac{1}{a b - a}\right) \cdot (\left(5 a b - 5 a\right) h_1 + \left(-a + b\right) h_2))\)\allowbreak\quad \((2,\left(\frac{1}{2 a b - 2 a}\right) \cdot (\left(-5 a b + 5 b\right) h_1 + \left(-a + b\right) h_3))\)\\
2 & \((1,\left(\frac{-1}{a^{2} b - a b^{2} - a^{2} + a b}\right) \cdot (\left(2 a^{2} b - 2 a b^{2} - 2 a^{2} + 2 a b\right) h_1 + \left(-a^{2} b + 2 a^{2} - 2 a b + 2 b^{2} - b\right) h_2 + \left(a^{2} b - a^{2} - a b + a\right) h_3))\)\allowbreak\quad \((2,\left(\frac{-a + 1}{2 a^{2} b - 2 a b^{2} - 2 a^{2} + 2 a b}\right) \cdot (\left(7 a b - 7 b^{2}\right) h_1 + \left(-b^{2} + b\right) h_2 + \left(b^{2} - a\right) h_3))\)\\
4 & \((1,\left(\frac{-1}{a - b}\right) \cdot (\left(a b - a\right) h_1 + \left(-5 a + 5 b\right) h_2))\)\allowbreak\quad \((2,\left(\frac{-1}{2 a - 2 b}\right) \cdot (\left(-a b + b\right) h_1 + \left(-a + b\right) h_3))\)\\
6 & \((1,\left(\frac{1}{a b - b^{2} - a + b}\right) \cdot (\left(a^{2} b - 2 a b^{2} + 2 a b - 2 a + b\right) h_1 + \left(-2 a b + 2 b^{2} + 2 a - 2 b\right) h_2 + \left(-a^{2} + a b + a - b\right) h_3))\)\allowbreak\quad \((2,\left(\frac{a - 1}{2 a^{2} b^{2} - 2 a b^{3} - 4 a^{2} b + 4 a b^{2} + 2 a^{2} - 2 a b}\right) \cdot (\left(-a^{2} b^{2} + 2 a b^{3} - a^{2} b - b^{3} + 2 a b - b^{2}\right) h_1 + \left(9 a b^{2} - 9 b^{3} - 9 a b + 9 b^{2}\right) h_2 + \left(a^{2} b - 2 a b^{2} + b^{3} + a^{2} - 2 a b + b^{2}\right) h_3))\)\\
7 & \((1,\left(\frac{-1}{a^{3} b - a^{2} b^{2} - a^{3} + a b^{2} + a^{2} - a b}\right) \cdot (\left(2 a^{3} b^{2} - a^{2} b^{3} - 3 a^{3} b + a^{2} b^{2} - a b^{3} + a^{3} - a^{2} b + 3 a b^{2} + a^{2} - 2 a b\right) h_1 + \left(-a^{3} b + a b^{3} - a^{3} + 7 a^{2} b - 7 a b^{2} + b^{3} - a^{2} + b^{2}\right) h_2 + \left(-9 a^{3} b + 9 a^{2} b^{2} + 9 a^{3} - 9 a b^{2} - 9 a^{2} + 9 a b\right) h_3))\)\allowbreak\quad \((2,\left(\frac{b}{2 a^{2} b - 2 a b^{2} - 2 a^{2} + 2 a b}\right) \cdot (\left(2 a^{2} b - a b^{2} - 2 a b - a + 2 b\right) h_1 + \left(-a b + b^{2} + a - b\right) h_2 + \left(-4 a^{2} + 4 a b + 4 a - 4 b\right) h_3))\)\\
9 & \((1,\left(\frac{-6}{a - b}\right) \cdot (\left(a b - a\right) h_1 + \left(-a + b\right) h_2))\)\allowbreak\quad \((2,\left(\frac{-1}{a - b}\right) \cdot (\left(-3 a b + 3 b\right) h_1 + \left(-a + b\right) h_3))\)\\
11 & \((2,\left(\frac{-a b + b}{a b - a}\right) \cdot p_{31})\)\\
12 & \((2,\left(\frac{a b - b}{a b - a}\right) \cdot p_{32})\)\\
13 & \((1,\left(\frac{a - b}{a b - a}\right) \cdot h_2)\)\allowbreak\quad \((2,\left(\frac{-a + b}{2 a b - 2 a}\right) \cdot h_3)\)\\
14 & \((1,\left(\frac{a - b}{a b - a}\right) \cdot p_{21})\)\\
15 & \((2,\left(\frac{-a + b}{2 a b - 2 a}\right) \cdot p_{31})\)\\
16 & \((1,\left(-1\right) \cdot p_{12})\)\allowbreak\quad \((2,\left(\frac{a b - b}{2 a b - 2 a}\right) \cdot p_{12})\)\\
17 & \((1,\left(-1\right) \cdot h_1)\)\allowbreak\quad \((2,\left(\frac{1}{2 a b - 2 a}\right) \cdot (\left(a b - b\right) h_1 + \left(-a + b\right) h_3))\)\\
18 & \((2,\left(\frac{-a + b}{2 a b - 2 a}\right) \cdot p_{32})\)\\
19 & \((1,\left(-1\right) \cdot p_{13})\)\allowbreak\quad \((2,\left(\frac{a b - b}{2 a b - 2 a}\right) \cdot p_{13})\)\\
20 & \((1,\left(\frac{-a + b}{a b - a}\right) \cdot p_{23})\)\\
21 & \((1,\left(\frac{1}{a b - a}\right) \cdot (\left(-a b + a\right) h_1 + \left(-a + b\right) h_2))\)\allowbreak\quad \((2,\left(\frac{a b - b}{2 a b - 2 a}\right) \cdot h_1)\)\\
\end{longtable}
\endgroup

Put $\Delta_a=a(b-1)$, $\Delta_b=b(a-1)$ and
$\Delta_{ab}=a-b$. Let $S_0,S_1$ select rows $(1,4)$ and $(13,17)$
of $M$, and put $M_0=S_0M$, $M_1=S_1M$. With
\[
 X_0=\Delta_bh_1h_2,\quad Y_0=\Delta_ah_1h_3,
 \quad Z_0=\Delta_{ab}h_2h_3,
\]
the determinant outputs are
\[
 \det M_0=\frac{2(6X_0+Y_0+Z_0)}{\Delta_a},\qquad
 \det M_1=\frac{\Delta_{ab}(X_0-Y_0-Z_0)}{2\Delta_a^2}.
\]
The following combined adjugate gives a left inverse without
choosing either determinant as a denominator.

\begin{equation}\label{supp:fnd:uleftinverse}
N=\frac{\Delta_a}{2}\operatorname{adj}(M_0)S_0+
\frac{2\Delta_a^2}{\Delta_{ab}}\operatorname{adj}(M_1)S_1
\end{equation}

\[
 NM=7\Delta_bh_1h_2\operatorname{Id}_2,\qquad
 \boldsymbol u=\frac{N\boldsymbol r}{7\Delta_bh_1h_2}.
\]
These are the outputs used in Proposition~\ref{fnd:finiteclosure}
and in the normalization of Section~\ref{sec:normalization}.

\clearpage
\section{Normalization and the lower ranks}\label{supp:normalization}

\subsection{The normalization factors and constant ratios}

The factors $\mu_i$ and the retained state $\mathsf X$ are those of
\eqref{fnd:mu} and \eqref{fnd:X}. The cubic reconstruction,
together with the left inverse above, removes
$\tau_4,\tau_5,\tau_0,u_0,u_1$ on
$h_1h_2h_3\mu_1\mu_2\mu_3\ne0$.
For constant abnormality ratios, use
$\alpha_{ij}=1-\lambda_j/\lambda_i$ and the three constants $R_i$
of Lemma~\ref{fnd:constantratio}. The matrix of the squared
off-diagonal coefficients is
\[
\begin{pmatrix}
R_1(R_1+1)&-\dfrac{\alpha_{31}}{\alpha_{21}}R_2&-\dfrac{\alpha_{13}}{\alpha_{23}}R_1R_3\\
-\dfrac{\alpha_{21}}{\alpha_{31}}R_1R_2&R_2(R_2+1)&-\dfrac{\alpha_{12}}{\alpha_{32}}R_3\\
-\dfrac{\alpha_{23}}{\alpha_{13}}R_1&-\dfrac{\alpha_{32}}{\alpha_{12}}R_2R_3&R_3(R_3+1)
\end{pmatrix}.
\]
Its determinant is $4$, using
$R_1R_2R_3=1$ and
$\alpha_{12}\alpha_{23}\alpha_{31}/(\alpha_{21}\alpha_{32}\alpha_{13})=-1$.

\subsection{The open normalization divisor}

On $\mu_1=0$, use the variables $q,d,g,z,X,Y,U,V,W$ of
Appendix~\ref{appn:mugeometry}. The opposite-pair equations are

\begin{align}
(X^2+d)Y+2z(1+z/d)XW&=0,\label{supp:fnd:mue2}\\
(dq-Y^2)X+2z(1-z/d)YW&=0,\label{supp:fnd:mue3}\\
2X^2Y^2+(d-z)Y^2-q(d+z)X^2&=0,\label{supp:fnd:mubiquad}\\
-(1+aq)(z^2+d^2)+2(1+qg)z-4XU-4gYV&=0.\label{supp:fnd:muquadratic}
\end{align}

\begin{equation}\label{supp:fnd:muuv}
UX=-\frac{gW^2}{ad}-\frac{(z-d)^2}{4d},\qquad
VY=\frac{W^2}{ad}+\frac{q(z+d)^2}{4d}.
\end{equation}

Set $R=X^2$ and $e=ad-a-1$. The two necessary cubics in $R$
have the short combination

\begin{equation}\label{supp:fnd:mu-short-cubic}
 F_1+F_2=4adz(z+d)H,
 \qquad H=4R(R+d)^2+eK,
\end{equation}

\[
 \begin{aligned}
 K={}&2(d^2+3dz-z^2)R^2\\
 &+(2d^3+4d^2z-9dz^2+z^3)R
       +3dz^2(z-d).
 \end{aligned}
\]

The Laurent roots at infinity and the first cubic evaluated on
those roots are

\begin{equation}\label{supp:fnd:mu-newton-roots}
 \begin{aligned}
 R_0&=-3d+O(z^{-1}),\\
 R_1&=\tfrac12z+\tfrac1{2e}+O(z^{-1}),\\
 R_2&=\tfrac12ez^2+O(z).
 \end{aligned}
\end{equation}

\begin{equation}\label{supp:fnd:mu-newton-values}
 \begin{aligned}
 F_1(R_0)&=12ade\,z^6+O(z^5),\\
 F_1(R_1)&=a(d-1)(ad-1)z^6+O(z^5),\\
 F_1(R_2)&=ae^3z^9+O(z^8).
 \end{aligned}
\end{equation}

\begin{equation}\label{supp:fnd:mu-short-resultant}
 \operatorname{Res}_R(F_1,H)
 =-768a^3d(d-1)(ad-1)e^4z^{21}+O(z^{20}).
\end{equation}

The eight cubic coefficients are given in
Appendix~\ref{appn:mucubics}. On the exceptional value $e=0$,
$d=1+a^{-1}$ and $H=4R(R+d)^2$. These outputs are used in
Lemma~\ref{fnd:muwall}.

\subsection{Two common similarity generators}

The orbit matrices and their constant eigenvalue equations are

\begin{equation}\label{supp:fnd:orbitlambda}
\det\{\Gamma-\lambda_i\operatorname{diag}(r^2,1)\}=0.
\end{equation}

\begin{equation}\label{supp:fnd:orbitcoeffs}
(a_0-\lambda_i)(n_{22}-\lambda_i)=0,\qquad
n_{11}(n_{22}-\lambda_i)-n_{12}^2=0.
\end{equation}

\begin{equation}\label{supp:fnd:twoplaneidentity}
(\lambda_j-\lambda_k)^2(U_j^2+V_j^2)
=(\lambda_j-m)(m-\lambda_k).
\end{equation}

For the nonabelian similarity pair, the harmonic height functions
and constant planar determinant satisfy

\begin{equation}\label{supp:fnd:nonabelianenergy}
\begin{gathered}
|F_\perp|^2+s=E,\qquad |F_\perp|^2/J_\perp^2+t/\kappa^2+1/\kappa^2=E',\\
t=c_0s+d_0,\quad c_0=\kappa^2/J_\perp^2>0.
\end{gathered}
\end{equation}

\begin{align}
b_0&=J_\perp c_0+J_\perp d_0/s,\label{supp:fnd:harmonicb}\\
a_0^2&=C_0-c_0s+C_1/s-J_\perp^2d_0^2/s^2,\label{supp:fnd:harmonicP}
\end{align}

For the commuting dilation and rotation, the two norm equations
reduce to the quadratic alternative

\begin{equation}\label{supp:fnd:quadraticnorm}
\mathcal N=u+b_{\mathrm n}x^2>0,\quad ub_{\mathrm n}\ne0,\qquad
y=\frac{Jx}{ku\sqrt{\mathcal N}},\qquad
\widetilde{\mathcal N}=\frac1u-\frac{ub_{\mathrm n}k^2}{J^2}y^2.
\end{equation}

The cylindrical phase uses $t=z/r$,
$Q=a_{\mathrm n}+mt^2$ and $\psi_c$ from
Appendix~\ref{appn:phaseproof}. The physical gradient and the
coefficient that removes the nonconstant phase are
\[
 F_f=\begin{pmatrix}
 a_{\mathrm n}/\sqrt Q&0&mt/\sqrt Q\\
 \sqrt Q(\omega-t\psi_c')&k\sqrt Q&\sqrt Q\psi_c'\\
 0&0&\nu
 \end{pmatrix},\qquad \det F_f=a_{\mathrm n}k\nu,
\]
\[
 [t^4]P=k^2m^2+(K_0-\omega m)^2+c_1^2.
\]
The remaining constant norm identity is

\begin{equation}\label{supp:fnd:constantnormpoly}
(E/a_{\mathrm n}-1)X_0W_0=\omega^2X_0W_0^2+\ell^2X_0^2+k^2W_0^2+h^2.
\end{equation}

\clearpage
\section{Common Euclidean isometries}\label{supp:isometries}

\subsection{Translation fibres}

For two translations, use the metric variables of
Appendix~\ref{appiso-tt}. The Hessian line, third derivative
system and nonzero quadratic coefficient are

\begin{equation}\label{supp:iso:tt-line}
 -2G_{nt}d\tau^2a_H+N b_H+G_{nt}KP\tau^4+G_{nt}Pd\tau^2-G_{nt}Qd_0\tau^4=0.
\end{equation}

\begin{equation}\label{supp:iso:tt-third}
 M=\begin{pmatrix}u&w\\-w&u\end{pmatrix},\quad
 u=-\frac{2d}{d_0\tau^2},\quad w=\frac{N}{G_{nt}d_0\tau^4},\quad
 \det M=\frac{4G_{nt}^2d^2\tau^4+N^2}{G_{nt}^2d_0^2\tau^8}>0.
\end{equation}

\begin{equation}\label{supp:iso:tt-leading}
 -\frac{G_{nt}\tau^6 H(\tau)}{d_0^2D(\tau)^3},\qquad
 D=-\tau^4G_{nt}^2,\qquad \deg H=24,\qquad [\tau^{24}]H=-K^6.
\end{equation}

\begin{equation}\label{supp:appiso-tt-adjugate}
 \widehat{\mathcal C}
 =\det M\,\mathcal C|_{c_H=e_H=0}
   -(\partial_{c_H}\mathcal C,\partial_{e_H}\mathcal C)
                   \operatorname{adj}(M)r.
\end{equation}

For a translation paired with a rotation or a screw motion, the
variables $v,\mathcal N_b,L,U,A,T,Q,D$ are those of
Appendix~\ref{appiso-trs-source}. The metric reconstruction and
the determinant outputs are

\begin{equation}\label{supp:iso:trs-shape}
 a_q=U,\quad a_\upsilon=QLU/v,\quad
 U^2=\frac{\mathcal N_b}{Q(v+L^2)},\quad
 P=v^3-\mathcal I_1v^2+\mathcal I_2v-j,
\end{equation}

\begin{equation}\label{supp:iso:trs-AT}
 A=\frac jv-\frac{P}{v\mathcal N_b},\quad
 R_0=\mathcal I_2-A-\frac{j(1+\mathcal N_b)}v,\quad
 T^2=vAR_0-jv^2.
\end{equation}

\begin{equation}\label{supp:appiso-G-reconstruct}
 G_{ij}=\frac AjK_{ij}+
       \frac v{\mathcal N_b}(s_0-2A/j)a_i a_j
       -\frac{T}{j\mathcal N_b}(a_i r_j+r_i a_j).
\end{equation}

Here $Q=2q+p^2$. The four coefficient rows, in the order phase
closedness, inverse vertical equilibrium, forward vertical
equilibrium and harmonicity, are as follows.

Subscripts $v,\mathcal N_b$
in what follows hold $\mathcal I_1,\mathcal I_2,j$ fixed.  Put $D=v+L^2$ and
\begin{align*}
 A_{\mathcal N_b}&=(j/v-A)/\mathcal N_b,\\
 T_{\mathcal N_b}&=\{v(R_0-A)A_{\mathcal N_b}-jA\}/(2T),\\
 A_v&=-j/v^2-(vP_v-P)/(v^2\mathcal N_b),\\
 T_v&=\{\mathcal I_2A-A^2-2jv+v(R_0-A)A_v\}/(2T),\\
 F&=A-TL/v,&G_0&=AL+T,\\
 S&=\frac{Q\{T(v-L^2)-v(R_0-A)L\}}{v^2D},\\
 R_1&=\frac{vR_0+AL^2+2TL}{vD}.
\end{align*}
Let
\[
 b_q^0=2\mathcal N_b/Q+2\mathcal N_b/D,\quad
 b_q=(2\mathcal N_b,0,2\mathcal N_b L/D,0),\quad b_\upsilon=(0,2\mathcal N_b,0,2\mathcal N_b L/D).
\]
For $y=(U_q/U,U_\upsilon/U,L_q,L_\upsilon)^{\top}$, the four rows of $M$ are
\begin{align*}
 M_1&=(QL/v,-1,Q/v,0),\\
 M_2&=(Q,L,0,1),\\
 M_3&=Q\{(F,0,F_L,0)+F_{\mathcal N_b}b_q\}
                  +(0,G_0,0,(G_0)_L)+(G_0)_{\mathcal N_b}b_\upsilon,\\
 M_4&=S_{\mathcal N_b}b_q+(0,0,S_L,(R_1)_L)+(R_1)_{\mathcal N_b}b_\upsilon.
\end{align*}
The equations are
\begin{equation}\label{supp:iso:trs-matrix}
 My+\boldsymbol r_0+\frac{K_{C_f}}{U}\boldsymbol r_C+\frac{K_{C_f^{-1}}}{U}\boldsymbol r_{\mathrm{inv}}
       +\frac1U \boldsymbol r_{-1}+K_{C_f}\boldsymbol r_{C0}+U\boldsymbol r_U=0,
\end{equation}
where
\begin{align*}
 \boldsymbol r_0&=(2p^2L/v^2,\,2,\,2F+Q(2F_v+F_{\mathcal N_b}b_q^0),\,
                           2S_v+S_{\mathcal N_b}b_q^0)^{\top},\\
 \boldsymbol r_C&=(0,0,-1,0)^{\top},&\boldsymbol r_{\mathrm{inv}}&=(0,1,0,0)^{\top},\\
 \boldsymbol r_{-1}&=(2p/v^2,0,0,0)^{\top},&\boldsymbol r_{C0}&=(0,0,0,p/v)^{\top},\\
 \boldsymbol r_U&=(0,0,0,-2pQF/v^2)^{\top}.
\end{align*}

\begin{equation}\label{supp:iso:trs-wall}
 \det M=-\frac{2\mathcal N_b^2Q^2}{vD}\,W,\quad
 W=\alpha L+\beta,
\end{equation}

\begin{equation}\label{supp:iso:wall-norm}
 [\mathcal N_b^2]N_U=\frac{j^2QK_{C_f}^2}{v^5},\qquad
 [\mathcal N_b^{3/2}]N_U=0,\qquad
 [\mathcal N_b]N_U\big|_{K_{C_f}=0}=-\frac{36j^2Q^2}{v^4}.
\end{equation}

\begin{equation}\label{supp:iso:trs-L-resultant}
 524288j^6v^{16}P^2(5v^3-3\mathcal I_1v^2+\mathcal I_2v+j)^3 G^{-1}(v),
\end{equation}

The two remaining compatibilities are
\eqref{iso:trs-frobenius}. The asymptotic
coefficients in \eqref{supp:iso:wall-norm} give the finite fibre
condition on its determinant zero set.

\subsection{Cylindrical linearization and two rotations}

At a cylinder with fixed pressure data, the five radial weights are
\[
 (-2,-1,-1,1,1).
\]
The quadratic obstruction uses $E,k,\omega,\Delta,\mathscr T,H_*,Z$
defined in Appendix~\ref{appiso-cylinder}. Its coefficient
identity is

\begin{equation}\label{supp:appiso-conic-wall}
 (E-1)\mathscr T+\Delta H_*=E\omega^2Z,\qquad
 H_*=(E-1)(k^2-1)+(E+1)\omega^2.
\end{equation}

The physical Hessian and third derivative systems for two rotations
are

\begin{equation}\label{supp:iso:rr-H}
 H_{a3}=\Omega_t^{\rm axial}\times F_f\partial_{x_a}-F_f(\Omega_s^{\rm axial}\times \partial_{x_a}),\qquad H_{ab}=\partial_a\partial_b f.
\end{equation}

\begin{equation}\label{supp:iso:rr-T}
 T_{ab3}=\Omega_t^{\rm axial}\times H_{ab}-H_{\Omega_s^{\rm axial}\times \partial_{x_a},b}-H_{a,\Omega_s^{\rm axial}\times \partial_{x_b}},
 \qquad T_{abc}=\partial_a\partial_b\partial_c f.
\end{equation}

\begin{equation}\label{supp:iso:rr-quadrics}
 (Q_0,Q_1,Q_2)^{\top}=(-M_1M_0^{-1}\mid I_3)b(H,\Omega_s^{\rm axial},\Omega_t^{\rm axial}),
\end{equation}

For the first twelve rows, with the determinant differentiated
instead of its logarithm, put
\[
 D_0=s^2-t^2,\quad
 L_0=t^2(D_0+p^2)-s^2q^2-D_0\gamma^2,\quad R_0=L_0+D_0^2,
\]
\[
 P_{10}=p^2L_0^2+q^2R_0^2+
 p^2q^2\{(s^2+t^2)^2(p^2+q^2)+2D_0^2(s^2+t^2)\}.
\]
The selected determinant is
\[
 \det M_0=128pqst(s^2-t^2)\gamma^5P_{10}.
\]
Here $P_{10}>0$ on the real chart $pqst\gamma(s^2-t^2)\ne0$.
The three determinant rows are multiplied by $st\gamma$ when
one uses $(st\gamma)^2d^2\log\det F_f$.
The three homogeneous compatibility conics are those in
\eqref{supp:iso:rr-quadrics}; their source rows and material resultant
are specified in Appendices~\ref{appiso-rr-matrix}
and~\ref{appiso-material}. The remaining coefficient on
$H_{11}^3$, after solving seven Hessian rows, is
\[
 -\frac{2(t^2p^2+s^2q^2)}{s^2t^2\gamma(s^2-t^2)}.
\]
For the transverse compactness argument, the homogeneous escape plane is

\begin{equation}\label{supp:iso:escape-plane}
 \boldsymbol p^{[\varepsilon w_3^Z]}=p\mathsf J_{13}+q\mathsf J_{23},\qquad \tau_0^{[\varepsilon w_3^Z]}=pq,\quad \boldsymbol\tau_{\perp0}^{[\varepsilon w_3^Z]}=0,\quad \boldsymbol u^{[\varepsilon w_3^Z]}=0,
 \qquad(p,q)\ne(0,0).
\end{equation}

\begin{equation}\label{supp:iso:normal-seven}
 s^2+\frac{b-a}{a(b-1)}t^2,\quad sh,\quad th,\quad h^2,
 \quad z\left(s+\frac{(a-b)q}{a(b-1)}t\right),\quad zh,\quad z^2.
\end{equation}

The source forcing and its comparison with the physical field
length are

\begin{align}
 A_1=A_2=B_{*1}=B_{*2}&=0,\label{supp:iso:forcing-kernel}\\
 2apq(bB_{*3}-A_3)-a(b-1)p^2(K_1+\sigma)
                         -(a-b)q^2(K_2+\sigma)&=0.
\end{align}

\begin{equation}\label{supp:iso:compact-comparison}
 \mathcal P\bigl(\boldsymbol{\mathcal G}^{(0)}-\boldsymbol{\mathcal G}^{(1)}\bigr)=O(|\varepsilon|(N+|\varepsilon|)).
\end{equation}

For two screws, the matched cylindrical carrier has radial weights
$(-2,-1,-1,-1,1,1)$. The pitch relation fixes the axial stretch as
$\gamma=q/p$; the fixed-pressure equations are given in
Appendix~\ref{appiso-ss}.

\clearpage
\section{Common dilations and rotations}\label{supp:dilations}

\subsection{Homothety and pressure matrices}

The seventeen homothety equations and the eight angular and pressure
observations are those of Appendix~\ref{adil:state-order}.
The projector pressure matrix and pure dilation determinant are

\begin{equation}\label{supp:dil:pressure-matrix}
 \mathsf P(\boldsymbol p^{[w_3^Z]})=
 \begin{pmatrix}
 p_{23}^{[w_3^Z]}-p_{32}^{[w_3^Z]}&p_{13}^{[w_3^Z]}&-p_{12}^{[w_3^Z]}\\
 -p_{23}^{[w_3^Z]}&p_{31}^{[w_3^Z]}-p_{13}^{[w_3^Z]}&p_{21}^{[w_3^Z]}\\
 p_{32}^{[w_3^Z]}&-p_{31}^{[w_3^Z]}&p_{12}^{[w_3^Z]}-p_{21}^{[w_3^Z]}
 \end{pmatrix}.
\end{equation}

\begin{equation}\label{supp:dil:pure-final-minor}
 A_0B_1-A_1B_0
   =\frac{2\kappa^2\ell s(z^2-z+1)^2}{z-1}\ne0.
\end{equation}

The eight-column angular minor away from principal directions is

\begin{equation}\label{supp:dil:eight-minor}
 \frac{(b-a)s(b-1)}{4a^2b^3}
 \{b(a-1)s^2t^2+a(b-1)s^2-(a-b)t^2\}.
\end{equation}

The normal equations at a principal direction have the following
coefficient outputs. The coordinates and row selections are those
of Appendix~\ref{adil:normal-estimate}.

\begin{equation}\label{supp:dil:infinity-plane}
 \begin{aligned}
 \boldsymbol p^{[\varepsilon w_3^Z]}&=p\mathsf J_{13}+q\mathsf J_{23},\qquad \tau_0^{[\varepsilon w_3^Z]}=pq,\\
 \tau_1^{[\varepsilon w_3^Z]}&=\cdots=\tau_5^{[\varepsilon w_3^Z]}=\boldsymbol u^{[\varepsilon w_3^Z]}=0,\qquad (p,q)\ne(0,0).
 \end{aligned}
\end{equation}

\begin{equation}\label{supp:dil:quadratic-normal}
 s^2+\frac{b-a}{a(b-1)}t^2,\quad sh,\quad th,\quad h^2,\quad
 z\left(s+\frac{(a-b)q}{a(b-1)}t\right),\quad zh,\quad z^2.
\end{equation}

\begin{equation}\label{supp:dil:quadratic-pivots}
\begin{array}{c|c}
 q\ne0,\ T_*\ne0&
 -\dfrac{49152(b-a)(b-1)^2a^2q^5T_*}{b^4(a-1)^4}\\[1mm]
 q=0&\dfrac{41472(b-a)(b-1)^4a^4}{b^5(a-1)^5}\\[1mm]
 T_*=0&\dfrac{2048q^{10}(7q^2-4b+3)}{b(q^2+1)^3}\\[1mm]
 T_*=0,\ 7q^2-4b+3=0&
 -\dfrac{49152q^9}{(q^2+1)^2(3q^2-1)}
\end{array}
\end{equation}

\begin{equation}\label{supp:dil:middle-four}
 (\mathcal P\boldsymbol{\mathcal G})_{(47,21,49,13)}=\delta M_4(O,p,q)u+\delta^3\boldsymbol c_4(p,q)
           +O(\delta^4+\delta^2|u|+|u|^2).
\end{equation}

\begin{equation}\label{supp:dil:middle-centre}
 u_2=\left(0,0,0,-\frac{2r_0p^2q(1+r_0^2)}{pr_0+q}\right).
\end{equation}

\begin{equation}\label{supp:dil:middle-obstruction}
 \frac{8pq(pr_0-q)(1+r_0^2)^2}{pr_0+q}=0,\qquad
 8r_0^2p^4(1+r_0^2)^2=0\quad\hbox{when }q=pr_0.
\end{equation}

\begin{equation}\label{supp:dil:resonant-five}
 (\mathcal P\boldsymbol{\mathcal G})_{(47,21,13,22,35)}=\delta M_5V+\delta^3\boldsymbol c_5+\boldsymbol r_{\mathrm{rem}},\qquad
 \det M_5\big|_{\Delta_{\mathrm{res}}=0}
      =\frac{4096r_0^6p^7(r_0^2+a)}a.
\end{equation}

\begin{equation}\label{supp:dil:resonant-centre}
            M_5(0,0,0,0,C_*)^{\top}+\boldsymbol c_5=0
\end{equation}

\begin{equation}\label{supp:dil:normal-rate}
                            N=O(|\varepsilon|).
\end{equation}

\subsection{The cylindrical carrier matrices}

Use the reference circle and the ordered twenty-three coordinates of
Appendix~\ref{adil:cylindrical-matrices}. The three ambient
linearizations $L_i$ are defined by differentiating
\eqref{dil:normalized-fields}; their first seventeen rows are

\begin{equation}\label{supp:dil:L-matrix-definition}
 (L_i)_{jk}=\partial_k\mathcal Q_i^j+d_jv_i\delta_{jk},\quad 1\le k\le17,
 \qquad (L_i)_{j,17+k}=d_j(\mathsf S^{[r]})_j\delta_{ik}.
\end{equation}

The last six rows are the derivatives of
$w_iw-\Gamma_i v$ and $-\Gamma_i a_0$, where $w=a_0\times v$.
For the nine free rows in \eqref{dil:kernel}, the kernel matrix
$N_{\mathrm{ker}}$ satisfies

\begin{equation}\label{supp:dil:kernel-identities}
 J_*N_{\mathrm{ker}}=0,\qquad J_*L_iN_{\mathrm{ker}}=0\quad(i=1,2,3),\qquad
                         (D\mathcal T)N_{\mathrm{ker}}=0.
\end{equation}

\begin{equation}\label{supp:dil:restricted-matrices}
 R_0=\mathcal P_f\Big(\sum_iv_iL_i\Big)N_{\mathrm{ker}},\qquad
 A_0=\mathcal P_f\Big(\sum_iw_iL_i\Big)N_{\mathrm{ker}}.
\end{equation}

Tables~\ref{supp:matrix-R_general} and~\ref{supp:matrix-A_general}
give the full restricted matrices in this free basis. We parametrize
the reference circle by
\[
 p=\frac{2t}{1+t^2},\qquad q=\frac{1-t^2}{1+t^2}.
\]
Here $t$ is only the rational parameter of this circle. The endpoint
bases are those specified after \eqref{dil:kernel}.

\begingroup\small
\setlength{\tabcolsep}{3pt}
\renewcommand{\arraystretch}{1.16}
\begin{longtable}{@{}r>{\raggedright\arraybackslash}p{0.91\linewidth}@{}}
\caption{The radial matrix $R_0$ on $pq\ne0$. An entry $(j,c)$ in row $i$ means that the $(i,j)$ coefficient is $c$. Unlisted entries are zero.}\label{supp:matrix-R_general}\\
\toprule Row & Nonzero column entries\\\midrule\endfirsthead
\toprule Row & Nonzero column entries\\\midrule\endhead
\bottomrule\endfoot
1 & \((1,-1)\)\allowbreak\quad \((4,\left(4\right) \cdot (t - 1) \cdot t \cdot (t + 1) \cdot (t^{2} + 1)^{-2})\)\allowbreak\quad \((6,\left(-1\right) \cdot (t - 1) \cdot (t + 1) \cdot (t^{2} + 1)^{-1})\)\\
2 & \((2,\left(\frac{1}{4}\right) \cdot t^{-2} \cdot (b - 1)^{-1} \cdot a^{-1} \cdot (t - 1)^{2} \cdot (t + 1)^{2} \cdot (t^{2} + 1)^{-2} \cdot (-a t^{4} + b t^{4} + 4 a b t^{2} - 2 a t^{2} - 2 b t^{2} - a + b))\)\allowbreak\quad \((3,\left(-\frac{1}{2}\right) \cdot t^{-1} \cdot (b - 1)^{-1} \cdot a^{-1} \cdot b \cdot (a - 1) \cdot (t - 1)^{3} \cdot (t + 1)^{3} \cdot (t^{2} + 1)^{-2})\)\allowbreak\quad \((5,\left(\frac{1}{2}\right) \cdot t^{-1} \cdot (t - 1) \cdot (t + 1))\)\allowbreak\quad \((8,\left(\frac{1}{4}\right) \cdot t^{-2} \cdot (b - 1)^{-1} \cdot a^{-1} \cdot (t - 1) \cdot (t + 1) \cdot (t^{2} + 1)^{-3} \cdot (-a t^{8} + b t^{8} + 8 a b t^{6} - 4 a t^{6} - 4 b t^{6} - 6 a t^{4} + 6 b t^{4} + 8 a b t^{2} - 4 a t^{2} - 4 b t^{2} - a + b))\)\allowbreak\quad \((9,\left(-\frac{1}{2}\right) \cdot t^{-1} \cdot (b - 1)^{-1} \cdot a^{-1} \cdot (t^{2} + 1)^{-3} \cdot (2 a b t^{8} - a t^{8} - b t^{8} - 4 a t^{6} + 4 b t^{6} + 12 a b t^{4} - 6 a t^{4} - 6 b t^{4} - 4 a t^{2} + 4 b t^{2} + 2 a b - a - b))\)\\
3 & \((2,\left(-1\right) \cdot t^{-1} \cdot (t - 1)^{3} \cdot (t + 1)^{3} \cdot (t^{2} + 1)^{-2})\)\allowbreak\quad \((3,\left(-1\right) \cdot (t^{2} + 1)^{-2} \cdot (3 t^{4} - 2 t^{2} + 3))\)\allowbreak\quad \((8,\left(-1\right) \cdot t^{-1} \cdot (t - 1)^{4} \cdot (t + 1)^{4} \cdot (t^{2} + 1)^{-3})\)\allowbreak\quad \((9,\left(-2\right) \cdot (t - 1)^{3} \cdot (t + 1)^{3} \cdot (t^{2} + 1)^{-3})\)\\
4 & \((4,\left(-1\right) \cdot (t^{2} + 1)^{-2} \cdot (t^{2} - 2 t - 1) \cdot (t^{2} + 2 t - 1))\)\allowbreak\quad \((6,\left(-2\right) \cdot t \cdot (t^{2} + 1)^{-1})\)\\
5 & \((2,\left(-1\right) \cdot t^{-1} \cdot (b - 1)^{-1} \cdot a^{-1} \cdot (t - 1) \cdot (t + 1) \cdot b \cdot (a - 1))\)\allowbreak\quad \((3,\left(-2\right) \cdot (b - 1)^{-1} \cdot a^{-1} \cdot b \cdot (a - 1))\)\allowbreak\quad \((5,1)\)\allowbreak\quad \((8,\left(-1\right) \cdot t^{-1} \cdot (b - 1)^{-1} \cdot a^{-1} \cdot b \cdot (a - 1) \cdot (t - 1)^{2} \cdot (t + 1)^{2} \cdot (t^{2} + 1)^{-1})\)\allowbreak\quad \((9,\left(-2\right) \cdot (b - 1)^{-1} \cdot a^{-1} \cdot (t - 1) \cdot (t + 1) \cdot b \cdot (a - 1) \cdot (t^{2} + 1)^{-1})\)\\
6 & \((1,\left(-1\right) \cdot (t - 1) \cdot (t + 1) \cdot (t^{2} + 1)^{-1})\)\allowbreak\quad \((4,\left(-2\right) \cdot t \cdot (t^{2} + 1)^{-3} \cdot (t^{2} - 4 t + 1) \cdot (t^{2} + 4 t + 1))\)\allowbreak\quad \((6,\left(-1\right) \cdot (t^{2} + 1)^{-2} \cdot (t^{4} + 10 t^{2} + 1))\)\\
7 & \((1,\left(-1\right) \cdot (t - 1)^{2} \cdot (t + 1)^{2} \cdot (t^{2} + 1)^{-2})\)\allowbreak\quad \((4,\left(-2\right) \cdot (t - 1) \cdot t \cdot (t + 1) \cdot (t^{2} + 1)^{-2})\)\allowbreak\quad \((7,-1)\)\\
8 & \((2,\left(-\frac{1}{4}\right) \cdot t^{-2} \cdot (b - 1)^{-1} \cdot a^{-1} \cdot (-a + b) \cdot (t - 1)^{3} \cdot (t + 1)^{3} \cdot (t^{2} + 1)^{-1})\)\allowbreak\quad \((3,\left(\frac{1}{2}\right) \cdot t^{-1} \cdot (b - 1)^{-1} \cdot a^{-1} \cdot (t^{2} + 1)^{-1} \cdot (a b t^{4} - b t^{4} + 2 a b t^{2} - 4 a t^{2} + 2 b t^{2} + a b - b))\)\allowbreak\quad \((5,\left(-\frac{1}{2}\right) \cdot t^{-1} \cdot (t^{2} + 1))\)\allowbreak\quad \((8,\left(-\frac{1}{4}\right) \cdot t^{-2} \cdot (b - 1)^{-1} \cdot a^{-1} \cdot (t^{2} + 1)^{-2} \cdot (-a t^{8} + b t^{8} + 4 a b t^{6} - 4 b t^{6} + 8 a b t^{4} - 14 a t^{4} + 6 b t^{4} + 4 a b t^{2} - 4 b t^{2} - a + b))\)\allowbreak\quad \((9,\left(\frac{1}{2}\right) \cdot t^{-1} \cdot (b - 1)^{-1} \cdot a^{-1} \cdot (t - 1) \cdot (t + 1) \cdot (t^{2} + 1)^{-2} \cdot (2 a b t^{4} - a t^{4} - b t^{4} + 4 a b t^{2} - 6 a t^{2} + 2 b t^{2} + 2 a b - a - b))\)\\
9 & \((2,\left(-2\right) \cdot t \cdot (t^{2} + 1)^{-1})\)\allowbreak\quad \((3,(t - 1) \cdot (t + 1) \cdot (t^{2} + 1)^{-1})\)\allowbreak\quad \((8,\left(-2\right) \cdot (t - 1) \cdot t \cdot (t + 1) \cdot (t^{2} + 1)^{-2})\)\allowbreak\quad \((9,(t - 1)^{2} \cdot (t + 1)^{2} \cdot (t^{2} + 1)^{-2})\)\\
\end{longtable}
\endgroup

\begingroup\small
\setlength{\tabcolsep}{3pt}
\renewcommand{\arraystretch}{1.16}
\begin{longtable}{@{}r>{\raggedright\arraybackslash}p{0.91\linewidth}@{}}
\caption{The angular matrix $A_0$ on $pq\ne0$. An entry $(j,c)$ in row $i$ means that the $(i,j)$ coefficient is $c$. Unlisted entries are zero.}\label{supp:matrix-A_general}\\
\toprule Row & Nonzero column entries\\\midrule\endfirsthead
\toprule Row & Nonzero column entries\\\midrule\endhead
\bottomrule\endfoot
1 & \((4,\left(-8\right) \cdot t^{2} \cdot (t^{2} + 1)^{-2})\)\allowbreak\quad \((6,\left(2\right) \cdot t \cdot (t^{2} + 1)^{-1})\)\allowbreak\quad \((7,\left(-2\right) \cdot (t - 1)^{-1} \cdot (t + 1)^{-1} \cdot t)\)\\
2 & \((2,\left(-\frac{1}{2}\right) \cdot t^{-1} \cdot (b - 1)^{-1} \cdot a^{-1} \cdot (t - 1) \cdot (t + 1) \cdot (t^{2} + 1)^{-2} \cdot (-a t^{4} + b t^{4} + 4 a b t^{2} - 2 a t^{2} - 2 b t^{2} - a + b))\)\allowbreak\quad \((3,\left(-1\right) \cdot (b - 1)^{-1} \cdot a^{-1} \cdot (t^{2} + 1)^{-2} \cdot (-a t^{4} + b t^{4} + 4 a b t^{2} - 2 a t^{2} - 2 b t^{2} - a + b))\)\allowbreak\quad \((8,\left(-\frac{1}{2}\right) \cdot t^{-1} \cdot (b - 1)^{-1} \cdot a^{-1} \cdot (t - 1)^{2} \cdot (t + 1)^{2} \cdot (t^{2} + 1)^{-3} \cdot (-a t^{4} + b t^{4} + 4 a b t^{2} - 2 a t^{2} - 2 b t^{2} - a + b))\)\allowbreak\quad \((9,\left(-1\right) \cdot (b - 1)^{-1} \cdot a^{-1} \cdot (t - 1) \cdot (t + 1) \cdot (t^{2} + 1)^{-3} \cdot (-a t^{4} + b t^{4} + 4 a b t^{2} - 2 a t^{2} - 2 b t^{2} - a + b))\)\\
3 & \((2,\left(2\right) \cdot (t - 1)^{2} \cdot (t + 1)^{2} \cdot (t^{2} + 1)^{-2})\)\allowbreak\quad \((3,\left(4\right) \cdot (t - 1) \cdot t \cdot (t + 1) \cdot (t^{2} + 1)^{-2})\)\allowbreak\quad \((8,\left(2\right) \cdot (t - 1)^{3} \cdot (t + 1)^{3} \cdot (t^{2} + 1)^{-3})\)\allowbreak\quad \((9,\left(4\right) \cdot t \cdot (t - 1)^{2} \cdot (t + 1)^{2} \cdot (t^{2} + 1)^{-3})\)\\
4 & \((4,\left(4\right) \cdot (t - 1) \cdot t \cdot (t + 1) \cdot (t^{2} + 1)^{-2})\)\allowbreak\quad \((6,\left(-1\right) \cdot (t - 1) \cdot (t + 1) \cdot (t^{2} + 1)^{-1})\)\allowbreak\quad \((7,1)\)\\
6 & \((4,\left(16\right) \cdot (t - 1) \cdot (t + 1) \cdot t^{2} \cdot (t^{2} + 1)^{-3})\)\allowbreak\quad \((6,\left(-4\right) \cdot (t - 1) \cdot t \cdot (t + 1) \cdot (t^{2} + 1)^{-2})\)\allowbreak\quad \((7,\left(4\right) \cdot t \cdot (t^{2} + 1)^{-1})\)\\
7 & \((1,\left(2\right) \cdot (t - 1) \cdot t \cdot (t + 1) \cdot (t^{2} + 1)^{-2})\)\allowbreak\quad \((4,\left(-1\right) \cdot (t - 1)^{2} \cdot (t + 1)^{2} \cdot (t^{2} + 1)^{-2})\)\\
8 & \((2,\left(\frac{1}{2}\right) \cdot t^{-1} \cdot (b - 1)^{-1} \cdot a^{-1} \cdot (-a + b) \cdot (t - 1)^{2} \cdot (t + 1)^{2} \cdot (t^{2} + 1)^{-1})\)\allowbreak\quad \((3,(b - 1)^{-1} \cdot a^{-1} \cdot (-a + b) \cdot (t - 1) \cdot (t + 1) \cdot (t^{2} + 1)^{-1})\)\allowbreak\quad \((8,\left(\frac{1}{2}\right) \cdot t^{-1} \cdot (b - 1)^{-1} \cdot a^{-1} \cdot (-a + b) \cdot (t - 1)^{3} \cdot (t + 1)^{3} \cdot (t^{2} + 1)^{-2})\)\allowbreak\quad \((9,(b - 1)^{-1} \cdot a^{-1} \cdot (-a + b) \cdot (t - 1)^{2} \cdot (t + 1)^{2} \cdot (t^{2} + 1)^{-2})\)\\
9 & \((2,\left(-1\right) \cdot (t - 1) \cdot (t + 1) \cdot (t^{2} + 1)^{-1})\)\allowbreak\quad \((3,\left(-2\right) \cdot t \cdot (t^{2} + 1)^{-1})\)\allowbreak\quad \((8,\left(-1\right) \cdot (t - 1)^{2} \cdot (t + 1)^{2} \cdot (t^{2} + 1)^{-2})\)\allowbreak\quad \((9,\left(-2\right) \cdot (t - 1) \cdot t \cdot (t + 1) \cdot (t^{2} + 1)^{-2})\)\\
\end{longtable}
\endgroup

At $(p,q)=(0,1)$, the endpoint matrices are

\begingroup\small
\setlength{\tabcolsep}{3pt}
\renewcommand{\arraystretch}{1.16}
\begin{longtable}{@{}r>{\raggedright\arraybackslash}p{0.91\linewidth}@{}}
\caption{The radial matrix at $(p,q)=(0,1)$. An entry $(j,c)$ in row $i$ means that the $(i,j)$ coefficient is $c$. Unlisted entries are zero.}\label{supp:matrix-R_pzero}\\
\toprule Row & Nonzero column entries\\\midrule\endfirsthead
\toprule Row & Nonzero column entries\\\midrule\endhead
\bottomrule\endfoot
1 & \((8,-1)\)\\
2 & \((2,-1)\)\allowbreak\quad \((6,1)\)\\
3 & \((3,(a - b)^{-1} \cdot (2 a b - 3 a + b))\)\allowbreak\quad \((5,\left(-2\right) \cdot (a - b)^{-1} \cdot (b - 1) \cdot a)\)\allowbreak\quad \((9,\left(-4\right) \cdot (a - b)^{-1} \cdot (b - 1) \cdot a)\)\\
4 & \((4,-1)\)\\
5 & \((3,\left(2\right) \cdot (a - b)^{-1} \cdot b \cdot (a - 1))\)\allowbreak\quad \((5,\left(-1\right) \cdot (a - b)^{-1} \cdot (2 a b - a - b))\)\allowbreak\quad \((9,\left(-4\right) \cdot (a - b)^{-1} \cdot b \cdot (a - 1))\)\\
6 & \((2,1)\)\allowbreak\quad \((6,-1)\)\\
7 & \((2,-1)\)\allowbreak\quad \((7,-1)\)\\
8 & \((1,-1)\)\\
9 & \((3,-1)\)\allowbreak\quad \((9,1)\)\\
\end{longtable}
\endgroup

\begingroup\small
\setlength{\tabcolsep}{3pt}
\renewcommand{\arraystretch}{1.16}
\begin{longtable}{@{}r>{\raggedright\arraybackslash}p{0.91\linewidth}@{}}
\caption{The angular matrix at $(p,q)=(0,1)$. An entry $(j,c)$ in row $i$ means that the $(i,j)$ coefficient is $c$. Unlisted entries are zero.}\label{supp:matrix-A_pzero}\\
\toprule Row & Nonzero column entries\\\midrule\endfirsthead
\toprule Row & Nonzero column entries\\\midrule\endhead
\bottomrule\endfoot
1 & \((3,-1)\)\allowbreak\quad \((5,1)\)\allowbreak\quad \((9,2)\)\\
4 & \((6,1)\)\allowbreak\quad \((7,1)\)\\
7 & \((4,-1)\)\\
8 & \((3,-1)\)\allowbreak\quad \((5,1)\)\allowbreak\quad \((9,2)\)\\
\end{longtable}
\endgroup

At $(p,q)=(1,0)$, they are

\begingroup\small
\setlength{\tabcolsep}{3pt}
\renewcommand{\arraystretch}{1.16}
\begin{longtable}{@{}r>{\raggedright\arraybackslash}p{0.91\linewidth}@{}}
\caption{The radial matrix at $(p,q)=(1,0)$. An entry $(j,c)$ in row $i$ means that the $(i,j)$ coefficient is $c$. Unlisted entries are zero.}\label{supp:matrix-R_qzero}\\
\toprule Row & Nonzero column entries\\\midrule\endfirsthead
\toprule Row & Nonzero column entries\\\midrule\endhead
\bottomrule\endfoot
1 & \((1,-1)\)\\
2 & \((9,-1)\)\\
3 & \((3,-1)\)\\
4 & \((4,1)\)\allowbreak\quad \((6,-1)\)\\
5 & \((3,\left(-2\right) \cdot (b - 1)^{-1} \cdot a^{-1} \cdot b \cdot (a - 1))\)\allowbreak\quad \((5,1)\)\\
6 & \((4,3)\)\allowbreak\quad \((6,-1 \cdot 3)\)\\
7 & \((4,1)\)\allowbreak\quad \((7,-1)\)\\
8 & \((3,1)\)\allowbreak\quad \((5,-1)\)\allowbreak\quad \((8,-1)\)\\
9 & \((2,-1)\)\\
\end{longtable}
\endgroup

\begingroup\small
\setlength{\tabcolsep}{3pt}
\renewcommand{\arraystretch}{1.16}
\begin{longtable}{@{}r>{\raggedright\arraybackslash}p{0.91\linewidth}@{}}
\caption{The angular matrix at $(p,q)=(1,0)$. An entry $(j,c)$ in row $i$ means that the $(i,j)$ coefficient is $c$. Unlisted entries are zero.}\label{supp:matrix-A_qzero}\\
\toprule Row & Nonzero column entries\\\midrule\endfirsthead
\toprule Row & Nonzero column entries\\\midrule\endhead
\bottomrule\endfoot
1 & \((4,-1 \cdot 2)\)\allowbreak\quad \((6,1)\)\allowbreak\quad \((7,1)\)\\
2 & \((3,-1)\)\\
7 & \((1,-1)\)\\
9 & \((3,-1)\)\\
\end{longtable}
\endgroup

For $pq\ne0$, the kernel basis and the entries of $R_0,A_0$ do not
require division by $b-ap^2-q^2$. They can therefore be specialized
at the material resonance $b=ap^2+q^2$ in the same nine free
coordinates.

The three determinants below correspond to the $14\times14$
submatrices of the constraint Jacobian $J_*$ selected in
\eqref{dil:carrier-minors}: the generic chart $pq\ne0$, the
endpoint $(p,q)=(0,1)$, and the endpoint $(p,q)=(1,0)$, respectively.

\begin{equation}\label{supp:dil:carrier-determinants}
 \frac{2p^4q^2(a-1)^2}{a(b-1)^2},\qquad
 \frac{2(a-1)^2}{a^2(b-1)^2},\qquad
 \frac{2(b-a)(a-1)^2}{a(b-1)^3}.
\end{equation}

Each determinant is nonzero on its chart for positive, pairwise
distinct principal values. Together with
\eqref{supp:dil:kernel-identities}, these minors give
$\operatorname{rank}J_*=14$ and $\dim\ker J_*=23-14=9$.
The restrictions of the radial and angular linearizations to this
kernel satisfy

\begin{align}\label{supp:dil:matrix-identities}
 \det(xI-R_0)&=x(x+2)(x-1)^2(x+1)^5,\\
 R_0(R_0-I)(R_0+2I)(R_0+I)^2&=0,\nonumber\\
 \quad [R_0,A_0]&=0,\nonumber\\
 A_0\mathcal P_{\mathrm{Hermite}}(R_0)&=0,\qquad \mathcal P_{\mathrm{Hermite}}(R_0)^2=\mathcal P_{\mathrm{Hermite}}(R_0),\nonumber\\
 \operatorname{tr}\mathcal P_{\mathrm{Hermite}}(R_0)&=3,\nonumber
\end{align}

where the Hermite interpolation polynomial is
\begin{equation}\label{supp:dil:hermite-projector}
                  \mathcal P_{\mathrm{Hermite}}(x)=1+\frac53x-\frac{x^2}{12}
                            -\frac76x^3-\frac5{12}x^4.
\end{equation}

The interpolation conditions
$\mathcal P_{\mathrm{Hermite}}(0)=\mathcal P_{\mathrm{Hermite}}(1)=1$
and $\mathcal P_{\mathrm{Hermite}}(-2)=\mathcal P_{\mathrm{Hermite}}(-1)
=\mathcal P_{\mathrm{Hermite}}'(-1)=0$, together with the annihilating
polynomial in \eqref{supp:dil:matrix-identities}, give
\[
 \operatorname{im}\mathcal P_{\mathrm{Hermite}}(R_0)
     =\ker R_0\oplus\ker(R_0-I).
\]
The image consists of the one dimensional centre, tangent to the
equilibrium circle, and the two dimensional eigenspace of eigenvalue
$1$. The double zero at $-1$ annihilates its generalized eigenspace.
The identity $A_0\mathcal P_{\mathrm{Hermite}}(R_0)=0$ gives vanishing
of the angular linearization on the image of the projection.

The additional coordinate $\epsilon$ has radial eigenvalue $1$
and zero angular linearization. On the ten dimensional tangent
space of the containing manifold in Lemma~\ref{dil:carrier},
the radial linearization therefore has characteristic polynomial
\begin{equation}\label{supp:dil:carrier-spectrum}
             \det(xI-D\mathcal R)=x(x+2)(x-1)^3(x+1)^5.
\end{equation}
Its centre, positive eigenspace and negative generalized eigenspace
have dimensions one, three and six, respectively, as used in
Lemma~\ref{dil:incoming}.

\clearpage
\section{The relative projective boundary}\label{supp:boundary}

\subsection{Ratio coordinates and the valuation regions}

Use $\kappa_0,q,\eta,\zeta$ from \eqref{r3:ratios}. In this
section $k=a\kappa_0$ and $m=1+k$ are the spectral constants of
Appendix~\ref{appr3:sources}. The affine retained coordinates
$\mathsf R_i$ are specified by

\begin{equation}\label{supp:r3:RRreconstruction}
\begin{aligned}
\widehat\tau_1={}&-\widehat p_{12}(\widehat p_{13}+\widehat p_{31})
 +\left(1-\frac{\eta \zeta }{\eta -\zeta }\right)\widehat p_{23}\\
 &-\frac{\eta \zeta }{\eta -\zeta }\widehat p_{32}-\frac{\eta \zeta }{2(\eta -\zeta )}\mathsf R_1,\\
\widehat\tau_2={}&-\frac{2\eta }{\kappa_0 (1+\eta )}\widehat p_{12}
 +\frac{2\eta ^2}{\kappa_0 (1+\eta )}\widehat p_{21}-\frac{\eta }{\kappa_0 (1+\eta )}\mathsf R_3,\\
\widehat\tau_3={}&\widehat p_{21}(\widehat p_{23}+\widehat p_{32})
 +\frac{\zeta -\eta +\eta \zeta }{\kappa_0 (1-\zeta )}\widehat p_{13}\\
 &+\frac{\zeta }{\kappa_0 (1-\zeta )}\widehat p_{31}+\frac{\zeta }{2\kappa_0 (1-\zeta )}\mathsf R_2.
\end{aligned}
\end{equation}

For a formal approach, put $n=\operatorname{ord}\tau$ and
$s_0=\min_j\operatorname{ord}v_j$, where $v_j=\tau^2\mathsf R_j$.
The local chart conditions and ideal outputs are

\begin{longtable}{@{}>{\raggedright\arraybackslash}p{0.46\linewidth}>{\raggedright\arraybackslash}p{0.46\linewidth}@{}}
\caption{Valuation regions of the weighted boundary.}\label{supp:tab:r3-valuations}\\
\toprule Region & Valuation condition\\\midrule\endfirsthead
\toprule Region & Valuation condition\\\midrule\endhead
Weighted boundary & $s_0=0$\\
Ordinary & $s_0\ge n$\\
Subcritical & $n/2<s_0<n$\\
Critical & $s_0=n/2$\\
Supercritical & $0<s_0<n/2$\\
\bottomrule\end{longtable}

\begin{longtable}{@{}>{\raggedright\arraybackslash}p{0.305\linewidth}>{\raggedright\arraybackslash}p{0.305\linewidth}>{\raggedright\arraybackslash}p{0.305\linewidth}@{}}

\caption{Local ideals on the primary and critical strata.}\label{supp:tab:r3-primary-ideals}\\
\toprule Chart & Local ideal & Stratum\\\midrule\endfirsthead
\toprule Chart & Local ideal & Stratum\\\midrule\endhead

Original weighted boundary $\boldsymbol v\ne0$ & $\tau\in I$ locally & All one-, two- and three-coordinate $\boldsymbol v$ supports.\\

Original weighted boundary nonzero $\overline p$ cone & $\tau\in I$ locally & All nonzero amplitudes, including vanishing sums.\\

Ordinary exceptional family & Boundary contained in the ordinary boundary family & All points on a dominant component.\\

Critical single $\mathsf R$ & $\tau\in\sqrt I$ locally & Eight roots of the auxiliary equations; all fifteen $\overline p$ supports.\\

Critical two $\mathsf R$ & $\tau\in\sqrt I$ locally & Rank three, rank two and rank one coefficient strata.\\

Critical three $\mathsf R$ & $\tau\in\sqrt I$ locally & All supports of paired amplitudes and hypersurfaces.\\

\bottomrule\end{longtable}

\begin{longtable}{@{}>{\raggedright\arraybackslash}p{0.305\linewidth}>{\raggedright\arraybackslash}p{0.305\linewidth}>{\raggedright\arraybackslash}p{0.305\linewidth}@{}}

\caption{Subcritical supports with one nonzero component.}\label{supp:tab:r3-single-supports}\\
\toprule Support & Result & Identity or remaining charts\\\midrule\endfirsthead
\toprule Support & Result & Identity or remaining charts\\\midrule\endhead

1,4 & $\tau^4\in I$ & Proposition~\ref{r3short:pure-axis}.\\

2,8 & $\tau\in I$ & Nine generators; selected determinants $-25qk^3/4032$ and $-544qk^3/10935$.\\

3,6,9,12 & $1\in I$ & Units $8x^2y^2$, $16y^2z^2$, $16x^2w^2$, $8z^2w^2$.\\

5 & $\tau\in I$ & Eight roots; the unit factors $1+2z_0$ and $9+20z_0$ are complementary.\\

7,13 & $\tau\in\sqrt I$ & Both the conic and its complement.\\

10 & $\tau\in\sqrt I$ & $w^2\ne\pm k$ and each hypersurface $w^2=\pm k$.\\

11,14 & $1\in I$ & The odd relation has unit coefficient of $xyw$ or $yzw$.\\

15 & $\tau\in\sqrt I$ & The conic and its resonance hypersurfaces.\\

\bottomrule\end{longtable}

The transverse coordinates and amplitude symbols in these tables
are the local coordinates of Appendix~\ref{r3:boundarysources}.
The quotient determinants on a pure-axis chart, with
$\overline R_3=1$ and $\varrho=\overline R_2$, are
\[
 D_9=29\varrho^2+49q,\quad D_{15}=q-19\varrho^2,\qquad
 -\frac{21\varrho q}{2D_9},\quad
 \frac{\varrho(\varrho^2+16q)}{4D_{15}}.
\]
At $D_9=0$, the complementary values are
$D_{15}=960q/29$ and $\varrho^2+16q=415q/29$.
At the inactive paired resonance, the residual square coefficient is
\[
 -\frac{16q}{\kappa_0}
       \left((\kappa_0+\tfrac12)^2+3\right).
\]
For the subcritical three-coordinate paired chart, the ideal is
\[
 J=(sW_\ell,s\kappa,\delta W_\ell,
           \delta W_\varrho,\delta\kappa,s^2),
\qquad \det C_J=-\frac{5q\varrho}{1152\kappa_0^5}.
\]
The exact coordinate graph and the relations defining this coefficient
matrix are in Appendix~\ref{r3:boundarysources}.

\subsection{Ordinary normal ideals}

Use the ordinary family \eqref{r3:incomingfamily} and the
normal coordinates

\begin{equation}\label{supp:appr3:ordinarynormals}
 \overline p=(1,n_0,n_1,n_2,n_3,\ell),\qquad
 \boldsymbol R=(\lambda_{\mathrm{in}}(v)\ell+n_\perp,u,v),\qquad
 \boldsymbol n=(n_0,n_1,n_2,n_3,n_\perp,\tau).
\end{equation}

Put $t=n_3$, $w=v+2$, and
$y=(n_0+n_3,n_1,n_2,\tau)$. The ordinary ideal and the ideal at
$\mu=m\ell^2-k=0$ are

\begin{equation}\label{supp:appr3ord:Jnine}
 J_9=(y,un_\perp,wt,n_\perp^2,n_\perp t,t^2).
\end{equation}

\begin{equation}\label{supp:appr3ord:Jten}
 J_{10}=(y,un_\perp,wt,n_\perp^2,n_\perp t,\mu t^2,t^3).
\end{equation}

At $\nu=m\ell^2+k=0$, use
$y_0=n_0+n_3$, $y_1=n_1+mt^2$,
$y_2=n_2+m\ell t^2$, $y_3=\tau$. The ideal and the shifted
quadratic coefficient are

\begin{equation}\label{supp:appr3ord:Jeleven}
 J_{11}=(y_0,y_1,y_2,y_3,un_\perp,wt,n_\perp^2,n_\perp t,
                   \nu t^2,ut^2,t^3).
\end{equation}

\begin{equation}\label{supp:appr3ord:shiftedsquare}
 g_i(\ell,u)=q_{i,tt}-mM_{i,n_1}-m\ell M_{i,n_2}.
\end{equation}

The normal coefficient matrices have unit maximal minors in the
local coefficient rings specified in
Appendix~\ref{appr3ord:complete}. Their generators include
$\tau$. At the material value $k+q=0$, the replacement base point
$(\eta,\zeta)=(1,-2)$ gives coefficient determinants
\[
 \ell^5\frac{64512}{25}\frac{q^4}{\kappa_0^2},\qquad
 -\frac{155648q^2}{135\kappa_0^3}.
\]

\subsection{Algebraicity on the ordinary family}

The induced differential equations and first integral are

\begin{equation}\label{supp:appr3:ordinaryconnection}
 D_\zeta v=\frac{(v-2\eta )(v+2)}{\Theta},\qquad D_\zeta \ell=0,\qquad
 D_\zeta (u/\beta)=\frac{v+1-\eta }{\Theta}\,(u/\beta).
\end{equation}

\[
 \mathcal J_{\mathrm{in}}=
 \frac{u^2}{\beta^2(v+2)(v-2\eta)},\qquad
 D_\eta\mathcal J_{\mathrm{in}}=D_\zeta\mathcal J_{\mathrm{in}}=0.
\]
For the normal minors at $u=0$, the two leading factors have the
identity
\[
 5(4\kappa_0\ell^2+5q)-4(5\kappa_0\ell^2+4q)=9q.
\]
At $v=2\eta$, the highest $\zeta$ coefficient of the relevant
normal factor is $-16mq(\eta+1)^5\ell^2$.
At the positive amplitude hypersurface, with $N=a\kappa_0^2$,
\[
 f_1=21622N-989(k+q),\quad f_2=14006N+377(k+q),
 \qquad 377f_1+989f_2=22003428N.
\]
The associated differential obstruction is

\begin{equation}\label{supp:r3:residue-identity}
 d\left(\frac{\zeta}{v+2}\right)
 =-\frac{(\eta+1)v}{(v-2\eta)(v+2)^2}\,dv.
\end{equation}

\begin{equation}\label{supp:r3:logarithmic-obstruction}
 D_\zeta\log(u/\beta)=-\frac1{2\eta}.
\end{equation}

\subsection{The additional ordinary hypersurface}

The local coordinates of Appendix~\ref{appr3tau:wall} are

\begin{equation}\label{supp:appr3tau:wallpoint}
 u=0,\qquad v=-2,\qquad
 \ell_0^2=d=-\frac{(a-b)\widehat h_3 ^2}{b(a-1)}\ne0.
\end{equation}

\begin{equation}\label{supp:appr3tau:coordinates}
 \begin{aligned}
 \overline p_{13}&=\rho_0-k_\perp+\kappa\ell\tau,& \overline p_{21}&=\rho_1,\\
 \overline p_{23}&=\rho_2,& \overline p_{31}&=k_\perp,\\
 n_\perp&=\boldsymbol R_1-\lambda_{\mathrm{in}}(v)\ell.
 \end{aligned}
\end{equation}

The entire linear normal module, its cokernel identity and the
relative equations are

\begin{equation}\label{supp:appr3tau:linearmodule}
 \begin{aligned}
 M_{j5}&=uA_j,\\
 M_{j4}-M_{j1}&=wB_j,\\
 M_{j6}+\kappa\ell M_{j1}&=uC_j+wD_j+gE_j,
 \end{aligned}
\end{equation}

\begin{equation}\label{supp:appr3tau:J}
 J=(un_\perp,\;w k_\perp,\;u\tau,\;w\tau,\;g\tau).
\end{equation}

\begin{equation}\label{supp:appr3tau:rangeidentity}
 \begin{aligned}
 \Delta M_j-cC_{j1}M_2-cC_{j2}M_{18}
              -d_AM_{j1}M_6=0 .
 \end{aligned}
\end{equation}

\begin{equation}\label{supp:appr3tau:Q}
 \begin{gathered}
 Q=(n_\perp^2,\;k_\perp n_\perp,\;n_\perp\tau,\;P_{\mathrm{hom}}(k_\perp,\tau)),\\
 P_{\mathrm{hom}}(k_\perp,\tau)=\tau^2P\left(\frac{k_\perp}{\ell_0\tau}\right).
 \end{gathered}
\end{equation}

\begin{equation}\label{supp:appr3tau:relativeform}
 \begin{aligned}
 J_i+\mathcal E_i^{\mathrm{rem}}&=0 &&(1\le i\le5),\\
 Q_j+\mathcal E_{5+j}^{\mathrm{rem}}&=0 &&(1\le j\le4),\\
 \mathcal E_i^{\mathrm{rem}}&\in(W)(n)^2+(n)^3.
 \end{aligned}
\end{equation}

\begin{equation}\label{supp:appr3tau:bounds}
 \begin{aligned}
 k_\perp&=\ell_0y_0\tau+O(\tau^2),& n_\perp&=O(\tau^2),\\
 (\ell-\ell_0,u,v+2)&=O(\tau^2),&P(y_0)&=0.
 \end{aligned}
\end{equation}

\begin{equation}\label{supp:appr3tau:determinant}
 O_{jc}=R_{8,\mathsf L_2}R_{j,c}
             -R_{j,\mathsf L_2}R_{8,c}.
\end{equation}

\begin{equation}\label{supp:appr3tau:norm}
 \mathcal N=f_0^2-P_1f_0f_1+P_0f_1^2.
\end{equation}

\begin{equation}\label{supp:appr3tau:laurent}
 \mathcal N(T,\pm T)
   =N_\pm(a,b)\epsilon^{-30}+O(\epsilon^{-29}).
\end{equation}

Here $P$ is the monic quadratic of
\eqref{appr3tau:Q}; the normal cubic coefficient matrix $O$
has the row order specified in \eqref{appr3tau:determinant}.
Its quadratic norm is
\[
 \operatorname{Nm}_{P}(f_0+f_1y)
       =f_0^2-P_1f_0f_1+P_0f_1^2,
 \qquad P(y)=y^2+P_1y+P_0.
\]
The two amplitude hypersurfaces have the determinant outputs

\begin{align}
 N_-&=\frac{9216}{49}
 \frac{a^2(b-1)^4(b-a)^{14}}
      {(a-1)^{18}b^{17}}\,C_-\notag\\
 &\qquad\times
 \frac{(7ab-a-6b)^2(11ab-3a-8b)^4}
      {(5ab-a-4b)^6(2ab-a-b)^2},
 \label{supp:appr3tau:negative}\\
 N_+&=\frac{9216}{49}
 \frac{a^2(b-1)^4(b-a)^{12}}
      {(a-1)^{18}b^{17}}\,C_+\notag\\
 &\qquad\times
 \frac{(5ab+a-6b)^2(5ab+3a-8b)^4}
      {(3ab+a-4b)^6},
 \label{supp:appr3tau:positive}
\end{align}

\clearpage
\section{Fibres over ratio curves}\label{supp:curves}

\subsection{The characteristic affine systems}

The curve slope is $\varsigma=d\zeta/d\eta$ and
$L=\eta\varsigma-\zeta$. The characteristic factors are
\[
 \varsigma-1,\quad L+1,\quad L+\varsigma,
 \quad\varsigma,\quad L,
\]
giving respectively $\zeta=\eta+c$, $\zeta=1+c\eta$,
$\zeta=c(\eta+1)$, $\zeta=c$, $\zeta=c\eta$.
For the first pencil, put $T=\zeta-\eta$ and
$s=\widehat p_{23}$; thus
$\widehat h_3=q\widehat h_2+qT/\kappa_0$.
On $T(T-1)s\ne0$, the three cubic
relations and $E_2L_{\mathrm p}=E_3L_{\mathrm p}=E_1B_{\mathrm p}=0$
have six-column determinant
\[
 \frac{4q^3T^2(1+\kappa_0\widehat h_2)
 (1-\kappa_0\widehat h_3/q)(1-T)^2s}{\kappa_0^2\widehat h_2}.
\]
The two equations $E_2^2L_{\mathrm p}=E_3^2L_{\mathrm p}=0$ have
$u_0,u_1$ coefficient determinant $-2q^2(1-T)^2$.
Here $L_{\mathrm p}=p_{33}-qp_{22}-(qT/\kappa_0)p_{11}$ and
$B_{\mathrm p}=p_{33}p_{23}+qp_{22}p_{32}$.
The affine $2\times2$ matrix and its Cramer numerators, including its
rank strata, are those of Appendix~\ref{appr3fib:pencils}.
On the common line $\zeta=\eta+1$, the reconstructed fibre has two
polynomials with coprime quintic leading forms; the two conics in
these forms are nonproportional, and one has determinant $-1$
(Appendix~\ref{r3:balancedline}).

\subsection{Directions, amplitudes and exceptional factors}

The opposite-pair direction coordinates are

\begin{equation}\label{supp:r3:directioncoordinates}
\begin{aligned}
 \widehat p_{12}&=w_3,&\widehat p_{13}&=w_2,&\widehat p_{23}&=w_1,\\
 \widehat p_{21}&=\eta w_3+\frac{((\eta +1)\varsigma -\zeta )m_3}{2\zeta },\\
 \widehat p_{31}&=-\zeta w_2-\frac{(1+\eta \varsigma -\zeta )m_2}{2\eta },\\
 \widehat p_{32}&=-\frac{\zeta w_1}{\eta }-\frac{(1-\varsigma )m_1}{2\eta },\\
 \mathsf R_1&=\frac{(\zeta -\eta \varsigma )m_1}{\eta \zeta },&
 \mathsf R_2&=\frac{\varsigma m_2}{\zeta },&\mathsf R_3&=-\frac{m_3}{\eta }.
\end{aligned}
\end{equation}

On nonzero pair support, the mixed equations give projective maps
between the directions. A nonidentity composition has at most two
fixed points. For a nonvertical projected factor $f(\eta,\zeta)$,
the contact slope is $-f_\eta/f_\zeta$.
The exceptional coefficient factors and their contact outputs are

\begin{longtable}{@{}>{\raggedright\arraybackslash}p{0.24\linewidth}>{\raggedright\arraybackslash}p{0.7\linewidth}@{}}
\caption{Coefficients on the exceptional direction hypersurfaces.}\label{supp:tab:r3-direction-coefficients}\\
\toprule Hypersurface & Nonzero coefficient or consequence\\\midrule\endfirsthead
\toprule Hypersurface & Nonzero coefficient or consequence\\\midrule\endhead
$L=0$ & $(\mathsf A_{\mathrm{amp}})_{31}=2(\eta-\zeta)$.\\
$\varsigma=0$ & $(\mathsf A_{\mathrm{amp}})_{12}=2(\zeta-1)$.\\
$L-\varsigma+1=0$ & The mixed coefficient is
$-4\eta^2((\eta-1)\varsigma+1)^2/[\varsigma(\varsigma-1)]$;
here $(\eta-1)\varsigma+1=\zeta$.\\
$L+\varsigma+1=0$ & The mixed coefficient is
$4\eta(\eta+\zeta)^2/(\zeta-1)$.
Its vanishing gives $\zeta=-\eta$, $\varsigma=-1$, and
$(\mathsf A_{\mathrm{amp}})_{31}=4\eta$.\\
$L+\varsigma-1=0$ & The mixed coefficient is
$-4\eta^4(\zeta+1)^2/(\eta-\zeta)$.
Its vanishing gives $\zeta=-1$, $\varsigma=0$, and
$(\mathsf A_{\mathrm{amp}})_{12}=-4$.\\
\bottomrule\end{longtable}

For two nonzero amplitudes $x,y$, the quartic combinations give
$\pm4x^2y^2$. With three amplitudes, the exceptional kernel lines
give $xy(x+y)^2$ and $xy(x^2+y^2)$.
The distinguished full-support direction is
\[
 t_1=\frac{\varsigma}{2\zeta},\qquad
 t_2=t_3=-\frac{L}{2\eta\zeta}.
\]
Its conics in the squared amplitudes have the coefficient minor
\[
 -2\left(\frac{\eta\varsigma}{\zeta}-1\right)^2.
\]
Their determinants and the two degenerate values are those of
Appendix~\ref{appr3fib:support}.

\subsection{The quadratic direction condition}

The quadratic factor is

\begin{equation}\label{supp:r3:quadratic-factor}
 \mathcal P_0=
 \eta\zeta(\eta^2-1)\varsigma^2
 -2(\eta^2-\zeta^2)\varsigma+\eta\zeta(1-\zeta^2).
\end{equation}

For a projected irreducible factor $f$, contact on
$\mathcal P_0=0$ has the exact form
\[
 \eta\zeta(\eta^2-1)f_\eta^2
 -2(\zeta^2-\eta^2)f_\eta f_\zeta
 -\eta\zeta(\zeta^2-1)f_\zeta^2=0.
\]
The leading and separation norm polynomials have total degrees $19$
and $20$, both of bidegree $(12,14)$; the rank-drop norm polynomial
has total degree nine. Their determinant and norm definitions are
given in Appendix~\ref{appr3fib:Bwall}.
For the distinguished direction,
$\varsigma'=\varsigma L/(\eta\zeta)$, and its contact condition is
\[
 \eta\zeta\bigl((\mathcal P_0)_\eta+
   \varsigma(\mathcal P_0)_\zeta\bigr)
   +\varsigma L(\mathcal P_0)_\varsigma=0.
\]
Writing the left side as $F_{\mathcal P_0}$, the resultant output is
\[
\begin{aligned}
 \operatorname{Res}_\varsigma(\mathcal P_0,F_{\mathcal P_0})
 ={}&-12(\zeta-\eta)(\zeta-1)(\zeta+1)(\eta-1)(\eta+1)\\
 &\cdot(\eta+\zeta)\zeta^5\eta^5(\eta^2+\zeta^2+1).
\end{aligned}
\]

\clearpage
\section{Moment matrices, norms and slope elimination}\label{supp:elimination}

\subsection{The two matrix presentations}

Set $K=1+\eta-\zeta$ and $c_{\mathrm m}=L/(2\eta\zeta)$.
The affine direction change and its quadratic coordinate are

\begin{equation}\label{supp:r3:momentcoordinates}
 t_1=\frac{V-\eta w_*}{L}+\frac{\varsigma }{2\zeta },\qquad
 t_2=\frac{U+w_*}{\varsigma }-\frac{L}{2\eta \zeta },\qquad
 t_3=w_*-\frac{L}{2\eta \zeta }.
\end{equation}

\begin{equation}\label{supp:r3:momentquadratics}
\begin{aligned}
 Q_*(w_*)&=2\zeta w_*\{2\eta \zeta Kw_*-(L^2+\eta \varsigma ^2-\zeta )\},\\
 Q(w)&=(2\eta \zeta w+L)\{2\zeta Kw-(\varsigma -1)(L+\varsigma )\}.
\end{aligned}
\end{equation}

The full matrices are determined by the polynomial matrix
$\mathsf A$ in Section~\ref{supp:momentcoefficients} and its five
row factors. They satisfy
\[
 M(w)=M_*(w+c_{\mathrm m})_{1:4,1:4},\qquad
 Q(w)=Q_*(w+c_{\mathrm m}).
\]
The coefficient extraction and reconstruction identities are

\begin{equation}\label{supp:r3:momentextraction}
 \operatorname{coeff}_{\xi}(F)=\left(
 \left.\partial_U\partial_VF\right|_{U=V=0},\,
 \left.\partial_UF\right|_{U=V=0},\,
 \left.\partial_VF\right|_{U=V=0},\,
 \frac{F(0,0,w)}{Q_*(w)}\right).
\end{equation}

\begin{equation}\label{supp:r3:momentreconstruction}
 F_i^{\mathrm{mom}}=\operatorname{coeff}_{\xi}(F_i^{\mathrm{mom}})\,(UV,U,V,Q_*)^{\top},
 \qquad (M_*)_{i,:}=\operatorname{coeff}_{\xi}(F_i^{\mathrm{mom}}),\quad 1\le i\le5.
\end{equation}

\subsection{The determinant and quadratic eliminants}

For a rational polynomial, $\operatorname{pp}$ denotes its primitive
polynomial numerator. Let $\operatorname{ng}_{g_*}$ remove only
the irreducible factors of
\[
 g_*=\eta\zeta(\eta-\zeta)(\zeta-1)(\eta+1)(1+\eta-\zeta)
        (\varsigma-1)(L+1)(L+\varsigma)\varsigma L
\]
and normalize the result as in Appendix~\ref{appr3:eliminants}.
If $\mathsf C$ is the $5\times3$ coefficient matrix of the five
maximal minors of $M_*$, then
\[
 \mathcal P_1=\operatorname{ng}_{g_*}
                 \gcd_{|I|=3}\det\mathsf C_I,
 \qquad (\deg\mathcal P_1,\deg_\varsigma\mathcal P_1)=(38,11).
\]
The residual minors $r_i$ have an identity
\[
 d(\eta)\varsigma^2(\varsigma-1)^2=\sum_{i=0}^5 A_ir_i,
 \qquad d(\eta)\ne0.
\]
For the second eliminant, put $D=\operatorname{pp}\det M$ and let
$a_j$ be column $j$ of $\operatorname{adj}M$. The quadratic kernel
equation gives
\[
 f_j=\operatorname{pp}\{Q(w)(a_j)_2(a_j)_3-(a_j)_1(a_j)_4\},
 \qquad R_j=\operatorname{Res}_w(D,f_j),\quad j=1,2.
\]
The common eliminant is
\[
 \mathcal P_2=\operatorname{ng}_{g_*}
       \bigl(\gcd(R_1,R_2)/\mathcal P_0\bigr),
\]
with
\[
 (\deg\mathcal P_2,\deg_\varsigma\mathcal P_2)=(48,14).
\]
The full coefficient representation of these outputs is specified by
the matrices in Section~\ref{supp:momentcoefficients}, the primitive
minor definitions and the fixed $(2,4)$ resultants in
Appendix~\ref{appr3:eliminants}.

\subsection{Schur reduction and the finite flat norm}

The determinant-one column transformation and transformed quadric are

\begin{equation}\label{supp:r3:Qcoefficients}
Q(w)=a_Qw^2+b_Qw+c_Q,\qquad
a_Q=4\eta \zeta ^2(1+\eta -\zeta ).
\end{equation}

\begin{equation}\label{supp:r3:columntransform}
P_*=
\begin{pmatrix}
1&\eta (w+b_Q/a_Q)&-w&\eta /a_Q\\
0&1&0&0\\
0&-\eta &1&0\\
0&-a_Qw&0&1
\end{pmatrix}.
\end{equation}

\begin{equation}\label{supp:r3:transformedmomentquadric}
\begin{aligned}
Qx_2x_3-x_1x_4
={}&w(a_Qy_1y_2+b_Qy_2y_3+y_3y_4)
  +c_Qy_2(y_3-\eta y_2)-y_1y_4\\
 &-\frac{\eta }{a_Q}y_4(b_Qy_2+y_4).
\end{aligned}
\end{equation}

On a constant-column minor chart, write
\[
 \widetilde{\mathsf N}=\begin{pmatrix}A_0&B_0(w)\\C_0&D_0(w)\end{pmatrix},
 \quad\Delta=\det A_0,\quad
 \mathsf R_{\mathrm{Schur}}=\Delta D_0-C_0\operatorname{adj}(A_0)B_0,
\]
\[
 \mathsf L_{\mathrm{Schur}}=
       \begin{pmatrix}-\operatorname{adj}(A_0)B_0\\\Delta\operatorname{Id}_2\end{pmatrix}.
\]
The exact block identities are

\begin{align}
 \det\mathsf R_{\mathrm{Schur}}&=\Delta\det\widetilde{\mathsf N},
 &\widetilde{\mathsf N}\mathsf L_{\mathrm{Schur}}&=\begin{pmatrix}0\\\mathsf R_{\mathrm{Schur}}\end{pmatrix},
 \label{supp:r3schur:basic}\\
 \mathsf L_{\mathrm{Schur}}\operatorname{adj}(\mathsf R_{\mathrm{Schur}})
 &=\Delta\operatorname{adj}(\widetilde{\mathsf N})_{\{3,4\}}.
 \label{supp:r3schur:columns}
\end{align}

\begin{align}
 \Delta^2\operatorname{adj}(\widetilde{\mathsf N})
 =\begin{pmatrix}
 (\det\mathsf R_{\mathrm{Schur}})E_0+E_0B_0J_0C_0E_0&-\Delta E_0B_0J_0\\
 -\Delta J_0C_0E_0&\Delta^2J_0
 \end{pmatrix}.
 \label{supp:r3schur:fulladj}
\end{align}

For $d(w)=d_2w^2+d_1w+d_0$ and
$f(w)=\sum_{i=0}^4f_iw^i$, put
\begin{align*}
 a_f={}&d_2^3f_1-d_1d_2^2f_2+d_2(d_1^2-d_0d_2)f_3
          +(-d_1^3+2d_0d_1d_2)f_4,\\
 b_f={}&d_2^3f_0-d_0d_2^2f_2+d_0d_1d_2f_3
          +(d_0^2d_2-d_0d_1^2)f_4.
\end{align*}
Thus $d_2^3f\equiv a_fw+b_f\pmod d$, and the norm output is

\begin{equation}
 \det\begin{pmatrix} b_f&-d_0 a_f\\
                      a_f&d_2 b_f-d_1 a_f\end{pmatrix}
 =d_2^3\operatorname{Res}_{2,4}(d,f).
 \label{supp:r3schur:normidentity}
\end{equation}

\subsection{Contact orders and elimination of the slope}

The differentiated curve conditions are

\begin{equation}\label{supp:r3:curve-differentiation}
 (\mathcal P_j)_\eta+\varsigma(\mathcal P_j)_\zeta
       +(\mathcal P_j)_\varsigma\varsigma'=0,
       \qquad j=1,2.
\end{equation}

Eliminating $\varsigma'$ gives
\[
 C_*=(\mathcal P_1)_\varsigma
       ((\mathcal P_2)_\eta+\varsigma(\mathcal P_2)_\zeta)
     -(\mathcal P_2)_\varsigma
       ((\mathcal P_1)_\eta+\varsigma(\mathcal P_1)_\zeta).
\]
The fixed Sylvester degrees of
$(\mathcal P_1,\mathcal P_2,C_*)$ are $(11,14,25)$. Put
$R=\operatorname{Res}_\varsigma(\mathcal P_1,\mathcal P_2)$,
$T=\operatorname{Res}_\varsigma(\mathcal P_1,C_*)$. The common divisor is
\[
\begin{aligned}
 D_0={}&\zeta^{12}(\zeta+1)^{20}(\zeta+\eta)^{20}
       (\zeta-\eta)^{25}(\zeta-1)^{25}\\
 &\cdot(\zeta-\eta-1)^{26}((\eta+1)\zeta-\eta)^6.
\end{aligned}
\]
The local orders are

\begin{longtable}{@{}>{\raggedright\arraybackslash}p{0.172\linewidth}>{\raggedright\arraybackslash}p{0.172\linewidth}>{\raggedright\arraybackslash}p{0.172\linewidth}>{\raggedright\arraybackslash}p{0.172\linewidth}>{\raggedright\arraybackslash}p{0.172\linewidth}@{}}

\caption{Intersection orders on the exceptional lines.}\label{supp:tab:r3-contact-orders}\\
\toprule Line in $\zeta $ & Order of $\mathcal P_1$ & Order of $\mathcal P_2$ & Lower order of $C_*$ & Common resultant order\\\midrule\endfirsthead
\toprule Line in $\zeta $ & Order of $\mathcal P_1$ & Order of $\mathcal P_2$ & Lower order of $C_*$ & Common resultant order\\\midrule\endhead

$\zeta =0$ & 4 & 3 & 6 & 12\\

$\zeta =-1$ & 4 & 5 & 7 & 20\\

$\zeta =-\eta $ & 4 & 5 & 7 & 20\\

$\zeta =\eta $ & 5 & 5 & 8 & 25\\

$\zeta =1$ & 5 & 5 & 8 & 25\\

$\zeta =\eta +1$ & 4 & 5 & 8 & 26\\

$\zeta =\eta /(\eta +1)$ & 2 & 3 & 4 & 6\\

\bottomrule\end{longtable}

The infinity calculation has coarse weight $487$ and a rank $23$
unit block. Its $2\times2$ Schur complement has determinant of order
at least six, giving $\deg_\zeta R\le481$.
At $p_*=32003$ and $\eta=3$, the exact finite-field outputs are
\[
 \deg_\zeta\overline R=481,\quad
 \deg_\zeta\overline D_0=134,\quad
 \deg_\zeta\overline{R/D_0}=347,
\]
\[
 \gcd(\overline{R/D_0},\overline{T/D_0})=1.
\]
The characteristic-zero result is
\[
 \gcd(R/D_0,T/D_0)=1\quad\hbox{in }\mathbb Q(\eta)[\zeta].
\]
The seven persistent ratio factors are

\begin{equation}\label{supp:r3:exceptional-ratio-loci}
 \zeta=0,\quad\zeta=-1,\quad\zeta=-\eta,\quad
 \zeta=\eta,\quad\zeta=1,\quad\zeta=\eta+1,\quad
 (\eta+1)\zeta=\eta.
\end{equation}

\clearpage
\section{Polynomial coefficients of the moment matrix}\label{supp:momentcoefficients}

The ratios $\eta,\zeta$, slope $\varsigma$ and direction parameter
$w$ are independent indeterminates in this coefficient list. They
are specialized to a solution curve after the matrix has been
formed. The definitions of $L,c_{\mathrm m}$ are those of
Section~\ref{supp:elimination}. The five row factors are

\[
\begin{aligned}
 d_{1}^{\mathrm{row}}&=\frac{421654016}{\varsigma \zeta ^2L},\qquad
 d_{2}^{\mathrm{row}}=d_{3}^{\mathrm{row}}=d_{4}^{\mathrm{row}}=\frac{6588344}{\varsigma \zeta ^2L},\\
 d_{5}^{\mathrm{row}}&=\frac{1}{2\eta \varsigma \zeta ^2(\eta +1)(\eta -\zeta )(\zeta -1)L}.
\end{aligned}
\]
The coefficient matrices satisfy
\[
\begin{aligned}
 M_*(w)&=\operatorname{diag}(d_{1}^{\mathrm{row}},\ldots,d_{5}^{\mathrm{row}})\mathsf A(w),\\
 M(w)&=\operatorname{diag}(d_{1}^{\mathrm{row}},\ldots,d_{4}^{\mathrm{row}})\mathsf A(w+c_{\mathrm m})_{1:4,\,1:4}.
\end{aligned}
\]
The following expressions define
$\mathsf A\in\operatorname{Mat}_{5\times4}(\mathbb Q[\eta ,\zeta ,\varsigma ,w])$.
The expressions $\gamma_i$ are defined successively in terms of
$\eta,\zeta,\varsigma,w$. They determine the twenty matrix entries.

\subsection{Auxiliary expressions}
\begingroup\small
\setlength{\abovedisplayskip}{5pt}
\setlength{\belowdisplayskip}{5pt}
\RenewDocumentEnvironment{aligned}{O{t}+b}{%
  \begingroup
  \let\aligned\suppaligned
  \let\endaligned\endsuppaligned
  \begin{array}[#1]{@{}l@{}}#2\end{array}%
  \endgroup
}{}


\[\begin{aligned}\gamma_{0}=\zeta ^{3}\\ \gamma_{1}=\eta  \gamma_{0}\\ \gamma_{2}=8 \gamma_{1}\end{aligned}\]

\[\begin{aligned}\gamma_{3}=- \zeta \\ \gamma_{4}=\eta  + \gamma_{3}\\ \gamma_{5}=\gamma_{4} + 1\end{aligned}\]

\[\begin{aligned}\gamma_{6}=\eta  - 1\\ \gamma_{7}=\gamma_{5} \gamma_{6}\\ \gamma_{8}=\eta  + \zeta \end{aligned}\]

\[\begin{aligned}\gamma_{9}=\zeta  + 1\\ \gamma_{10}=\eta  \zeta \\ \gamma_{11}=- \eta  + \zeta  + \gamma_{10}\end{aligned}\]

\[\begin{aligned}\gamma_{12}=\eta ^{2}\\ \gamma_{13}=\zeta ^{2}\\ \gamma_{14}=\eta  \gamma_{13}\end{aligned}\]

\[\begin{aligned}\gamma_{15}=\zeta  \gamma_{12}\\ \gamma_{16}=- \gamma_{15}\\ \gamma_{17}=2 \gamma_{10}\end{aligned}\]

\[\begin{aligned}\gamma_{18}=- \gamma_{12} - \gamma_{13} - \gamma_{14} - \gamma_{16} - \gamma_{17} - \gamma_{4}\\ \gamma_{19}=\gamma_{11} \gamma_{18} \gamma_{8} \gamma_{9}\\ \gamma_{20}=\gamma_{19} \gamma_{2} \gamma_{7}\end{aligned}\]

\[\begin{aligned}\gamma_{21}=\eta ^{3}\\ \gamma_{22}=- \gamma_{14}\\ \gamma_{23}=\gamma_{0} \gamma_{12}\end{aligned}\]

\[\begin{aligned}\gamma_{24}=34 \gamma_{23}\\ \gamma_{25}=\zeta ^{6}\\ \gamma_{26}=\gamma_{12} \gamma_{25}\end{aligned}\]

\[\begin{aligned}\gamma_{27}=\zeta ^{5}\\ \gamma_{28}=2 \gamma_{27}\\ \gamma_{29}=\varsigma  \gamma_{21}\end{aligned}\]

\[\begin{aligned}\gamma_{30}=- \gamma_{28} \gamma_{29}\\ \gamma_{31}=\zeta ^{4}\\ \gamma_{32}=- 18 \eta  \gamma_{31}\end{aligned}\]

\[\begin{aligned}\gamma_{33}=\eta ^{5}\\ \gamma_{34}=\gamma_{0} \gamma_{33}\\ \gamma_{35}=\eta ^{4}\end{aligned}\]

\[\begin{aligned}\gamma_{36}=\gamma_{12} \gamma_{27}\\ \gamma_{37}=- 9 \gamma_{36}\\ \gamma_{38}=\gamma_{31} \gamma_{35}\end{aligned}\]

\[\begin{aligned}\gamma_{39}=7 \gamma_{38}\\ \gamma_{40}=\gamma_{21} \gamma_{27}\\ \gamma_{41}=3 \gamma_{40}\end{aligned}\]

\[\begin{aligned}\gamma_{42}=\eta ^{6}\\ \gamma_{43}=2 \varsigma \\ \gamma_{44}=\gamma_{42} \gamma_{43}\end{aligned}\]

\[\begin{aligned}\gamma_{45}=\zeta  \gamma_{33}\\ \gamma_{46}=\eta  \gamma_{27}\\ \gamma_{47}=\zeta  \gamma_{35}\end{aligned}\]

\[\begin{aligned}\gamma_{48}=\gamma_{0} \gamma_{35}\\ \gamma_{49}=\gamma_{0} \gamma_{21}\\ \gamma_{50}=\gamma_{12} \gamma_{13}\end{aligned}\]

\[\begin{aligned}\gamma_{51}=21 \varsigma  \gamma_{50}\\ \gamma_{52}=- \gamma_{51}\\ \gamma_{53}=\gamma_{14} \gamma_{43}\end{aligned}\]

\[\begin{aligned}\gamma_{54}=- 2 \varsigma  \gamma_{13} \gamma_{42}\\ \gamma_{55}=2 \gamma_{31}\\ \gamma_{56}=\gamma_{29} \gamma_{55}\end{aligned}\]

\[\begin{aligned}\gamma_{57}=\eta  \varsigma \\ \gamma_{58}=\gamma_{31} \gamma_{57}\\ \gamma_{59}=4 \gamma_{58}\end{aligned}\]

\[\begin{aligned}\gamma_{60}=\varsigma  \gamma_{42}\\ \gamma_{61}=4 \zeta  \gamma_{60}\\ \gamma_{62}=\varsigma  \gamma_{31}\end{aligned}\]

\[\begin{aligned}\gamma_{63}=\gamma_{35} \gamma_{62}\\ \gamma_{64}=8 \gamma_{0}\\ \gamma_{65}=\varsigma  \gamma_{34}\end{aligned}\]

\[\begin{aligned}\gamma_{66}=\varsigma  \gamma_{23}\\ \gamma_{67}=\varsigma  \gamma_{35}\\ \gamma_{68}=\gamma_{0} \gamma_{67}\end{aligned}\]

\[\begin{aligned}\gamma_{69}=\gamma_{0} \gamma_{29}\\ \gamma_{70}=\zeta  \gamma_{21}\\ \gamma_{71}=\varsigma  \gamma_{70}\end{aligned}\]

\[\begin{aligned}\gamma_{72}=\varsigma  \gamma_{45}\\ \gamma_{73}=\zeta  \gamma_{67}\\ \gamma_{74}=\varsigma  \gamma_{15}\end{aligned}\]

\[\begin{aligned}\gamma_{75}=4 \gamma_{74}\\ \gamma_{76}=- \gamma_{27}\\ \gamma_{77}=\varsigma  \gamma_{12}\end{aligned}\]

\[\begin{aligned}\gamma_{78}=\gamma_{55} \gamma_{77}\\ \gamma_{79}=\gamma_{76} + \gamma_{78}\\ \gamma_{80}=\gamma_{0} + \gamma_{75} + \gamma_{79}\end{aligned}\]

\[\begin{aligned}\gamma_{81}=- \gamma_{31}\\ \gamma_{82}=2 \gamma_{25}\\ \gamma_{83}=\eta  \gamma_{82}\end{aligned}\]

\[\begin{aligned}\gamma_{84}=7 \gamma_{21}\\ \gamma_{85}=\gamma_{31} \gamma_{84}\\ \gamma_{86}=\gamma_{28} \gamma_{57}\end{aligned}\]

\[\begin{aligned}\gamma_{87}=\gamma_{81} + \gamma_{83} + \gamma_{85} - \gamma_{86}\\ \gamma_{88}=\varsigma  \gamma_{36}\\ \gamma_{89}=- 4 \gamma_{88}\end{aligned}\]

\[\begin{aligned}\gamma_{90}=7 \gamma_{70}\\ \gamma_{91}=\gamma_{25} + \gamma_{89} + \gamma_{90}\\ \gamma_{92}=4 \gamma_{0}\end{aligned}\]

\[\begin{aligned}\gamma_{93}=w^{2}\\ \gamma_{94}=\gamma_{13} \gamma_{84}\\ \gamma_{95}=2 \gamma_{14}\end{aligned}\]

\[\begin{aligned}\gamma_{96}=- \gamma_{15} \gamma_{43}\\ \gamma_{97}=\gamma_{0} \gamma_{84}\\ \gamma_{98}=\gamma_{57} \gamma_{92}\end{aligned}\]

\[\begin{aligned}\gamma_{99}=- \gamma_{98}\\ \gamma_{100}=\varsigma ^{2}\\ \gamma_{101}=\gamma_{100} \gamma_{33}\end{aligned}\]

\[\begin{aligned}\gamma_{102}=15 \gamma_{67}\\ \gamma_{103}=23 \gamma_{50}\\ \gamma_{104}=- 4 \varsigma  \gamma_{12} \gamma_{13}\end{aligned}\]

\[\begin{aligned}\gamma_{105}=\gamma_{100} \gamma_{50}\\ \gamma_{106}=2 \gamma_{105}\\ \gamma_{107}=\gamma_{100} \gamma_{70}\end{aligned}\]

\[\begin{aligned}\gamma_{108}=\gamma_{100} \gamma_{35}\\ \gamma_{109}=4 \gamma_{108}\\ \gamma_{110}=11 \gamma_{69}\end{aligned}\]

\[\begin{aligned}\gamma_{111}=\gamma_{13} \gamma_{29}\\ \gamma_{112}=28 \gamma_{111}\\ \gamma_{113}=2 \gamma_{13}\end{aligned}\]

\[\begin{aligned}\gamma_{114}=\gamma_{100} \gamma_{21}\\ \gamma_{115}=- \gamma_{113} \gamma_{114}\\ \gamma_{116}=\gamma_{115} + \gamma_{29}\end{aligned}\]

\[\begin{aligned}\gamma_{117}=\varsigma  \gamma_{14}\\ \gamma_{118}=\gamma_{27} - \gamma_{78}\\ \gamma_{119}=\gamma_{117} + \gamma_{118}\end{aligned}\]

\[\begin{aligned}\gamma_{120}=\eta  \gamma_{55}\\ \gamma_{121}=- \gamma_{120}\\ \gamma_{122}=\gamma_{33} \gamma_{43}\end{aligned}\]

\[\begin{aligned}\gamma_{123}=\gamma_{121} + \gamma_{122}\\ \gamma_{124}=- \gamma_{0}\\ \gamma_{125}=- \gamma_{85}\end{aligned}\]

\[\begin{aligned}\gamma_{126}=\eta  \gamma_{28}\\ \gamma_{127}=\varsigma  \gamma_{33}\\ \gamma_{128}=\gamma_{113} \gamma_{127} + \gamma_{124} + \gamma_{125} + \gamma_{126} + \gamma_{36}\end{aligned}\]

\[\begin{aligned}\gamma_{129}=2 \gamma_{1}\\ \gamma_{130}=- \gamma_{129}\\ \gamma_{131}=\gamma_{100} \gamma_{42}\end{aligned}\]

\[\begin{aligned}\gamma_{132}=2 \gamma_{131}\\ \gamma_{133}=\gamma_{130} + \gamma_{132}\\ \gamma_{134}=8 \gamma_{108}\end{aligned}\]

\[\begin{aligned}\gamma_{135}=- \gamma_{90}\\ \gamma_{136}=\gamma_{134} + \gamma_{135} + \gamma_{16}\\ \gamma_{137}=8 \gamma_{101}\end{aligned}\]

\[\begin{aligned}\gamma_{138}=\gamma_{13} \gamma_{67}\\ \gamma_{139}=\gamma_{13} \gamma_{137} + 23 \gamma_{138}\\ \gamma_{140}=\zeta  - 1\end{aligned}\]

\[\begin{aligned}\gamma_{141}=\gamma_{3} + \gamma_{57}\\ \gamma_{142}=\gamma_{113} \gamma_{140} \gamma_{141}\\ \gamma_{143}=- \gamma_{55}\end{aligned}\]

\[\begin{aligned}\gamma_{144}=\zeta ^{7}\\ \gamma_{145}=\eta ^{7}\\ \gamma_{146}=\gamma_{100} \gamma_{145}\end{aligned}\]

\[\begin{aligned}\gamma_{147}=\gamma_{113} \gamma_{33}\\ \gamma_{148}=3 \gamma_{1}\\ \gamma_{149}=\eta  \gamma_{25}\end{aligned}\]

\[\begin{aligned}\gamma_{150}=3 \gamma_{149}\\ \gamma_{151}=\gamma_{21} \gamma_{25}\\ \gamma_{152}=\zeta  \gamma_{146}\end{aligned}\]

\[\begin{aligned}\gamma_{153}=\gamma_{100} \gamma_{36}\\ \gamma_{154}=\varsigma  \gamma_{0}\\ \gamma_{155}=\gamma_{154} \gamma_{42}\end{aligned}\]

\[\begin{aligned}\gamma_{156}=2 \gamma_{155}\\ \gamma_{157}=\gamma_{43} \gamma_{70}\\ \gamma_{158}=2 \zeta \end{aligned}\]

\[\begin{aligned}\gamma_{159}=\varsigma  \gamma_{158}\\ \gamma_{160}=\gamma_{159} \gamma_{42}\\ \gamma_{161}=\gamma_{108} \gamma_{27}\end{aligned}\]

\[\begin{aligned}\gamma_{162}=7 \gamma_{108}\\ \gamma_{163}=7 \gamma_{0}\\ \gamma_{164}=\gamma_{13} \gamma_{21}\end{aligned}\]

\[\begin{aligned}\gamma_{165}=\gamma_{100} \gamma_{164}\\ \gamma_{166}=\zeta  \gamma_{131}\\ \gamma_{167}=\gamma_{12} \gamma_{62}\end{aligned}\]

\[\begin{aligned}\gamma_{168}=\gamma_{100} \gamma_{34}\\ \gamma_{169}=2 \gamma_{107}\\ \gamma_{170}=- \gamma_{169}\end{aligned}\]

\[\begin{aligned}\gamma_{171}=2 \gamma_{0}\\ \gamma_{172}=\gamma_{170} + \gamma_{171} \gamma_{57}\\ \gamma_{173}=- \gamma_{131}\end{aligned}\]

\[\begin{aligned}\gamma_{174}=- \gamma_{100} \gamma_{97} + \gamma_{105} + \gamma_{173}\\ \gamma_{175}=12 \gamma_{72}\\ \gamma_{176}=\gamma_{175} + \gamma_{28} + \gamma_{86}\end{aligned}\]

\[\begin{aligned}\gamma_{177}=\gamma_{13} \gamma_{146}\\ \gamma_{178}=3 \gamma_{114}\\ \gamma_{179}=\gamma_{100} \gamma_{45}\end{aligned}\]

\[\begin{aligned}\gamma_{180}=21 \gamma_{108}\\ \gamma_{181}=\gamma_{0} \gamma_{180} + \gamma_{177} + \gamma_{178} \gamma_{27} + 18 \gamma_{179}\\ \gamma_{182}=4 \gamma_{14}\end{aligned}\]

\[\begin{aligned}\gamma_{183}=\gamma_{182} w\\ \gamma_{184}=2 \gamma_{23}\\ \gamma_{185}=- \gamma_{184}\end{aligned}\]

\[\begin{aligned}\gamma_{186}=\gamma_{108} \gamma_{113}\\ \gamma_{187}=8 \gamma_{31}\\ \gamma_{188}=10 \gamma_{0}\end{aligned}\]

\[\begin{aligned}\gamma_{189}=4 \gamma_{50}\\ \gamma_{190}=\gamma_{12} \gamma_{55}\\ \gamma_{191}=\varsigma  \gamma_{113}\end{aligned}\]

\[\begin{aligned}\gamma_{192}=\gamma_{100} \gamma_{12}\\ \gamma_{193}=\eta  \gamma_{31}\\ \gamma_{194}=\gamma_{100} \gamma_{15}\end{aligned}\]

\[\begin{aligned}\gamma_{195}=\gamma_{0} \gamma_{57}\\ \gamma_{196}=\gamma_{100} \gamma_{17}\\ \gamma_{197}=\gamma_{158} \gamma_{67}\end{aligned}\]

\[\begin{aligned}\gamma_{198}=\zeta  \gamma_{109}\\ \gamma_{199}=4 \gamma_{105}\\ \gamma_{200}=17 \gamma_{117}\end{aligned}\]

\[\begin{aligned}\gamma_{201}=\gamma_{55} \gamma_{57}\\ \gamma_{202}=2 \gamma_{77}\\ \gamma_{203}=\gamma_{201} - \gamma_{202}\end{aligned}\]

\[\begin{aligned}\gamma_{204}=\gamma_{191} \gamma_{4}\\ \gamma_{205}=- \gamma_{1}\\ \gamma_{206}=- \gamma_{108}\end{aligned}\]

\[\begin{aligned}\gamma_{207}=7 \gamma_{12}\\ \gamma_{208}=\gamma_{0} \gamma_{207}\\ \gamma_{209}=\gamma_{100} \gamma_{208}\end{aligned}\]

\[\begin{aligned}\gamma_{210}=7 \gamma_{168}\\ \gamma_{211}=- \gamma_{210}\\ \gamma_{212}=- \gamma_{45}\end{aligned}\]

\[\begin{aligned}\gamma_{213}=\gamma_{31} \gamma_{33}\\ \gamma_{214}=\gamma_{13} \gamma_{33}\\ \gamma_{215}=\gamma_{27} \gamma_{84}\end{aligned}\]

\[\begin{aligned}\gamma_{216}=\gamma_{13} \gamma_{35}\\ \gamma_{217}=\gamma_{21} \gamma_{31}\\ \gamma_{218}=\gamma_{12} \gamma_{31}\end{aligned}\]

\[\begin{aligned}\gamma_{219}=\gamma_{1} \gamma_{100}\\ \gamma_{220}=\gamma_{29} \gamma_{31}\\ \gamma_{221}=\gamma_{108} \gamma_{13}\end{aligned}\]

\[\begin{aligned}\gamma_{222}=\gamma_{114} \gamma_{31}\\ \gamma_{223}=\gamma_{100} \gamma_{218}\\ \gamma_{224}=\gamma_{108} \gamma_{31}\end{aligned}\]

\[\begin{aligned}\gamma_{225}=10 \gamma_{67}\\ \gamma_{226}=\gamma_{101} \gamma_{13}\\ \gamma_{227}=4 \gamma_{25}\end{aligned}\]

\[\begin{aligned}\gamma_{228}=\varsigma  \gamma_{13}\\ \gamma_{229}=\gamma_{228} \gamma_{42}\\ \gamma_{230}=\gamma_{13} \gamma_{131}\end{aligned}\]

\[\begin{aligned}\gamma_{231}=10 \gamma_{13}\\ \gamma_{232}=\zeta  \gamma_{108}\\ \gamma_{233}=12 \gamma_{67}\end{aligned}\]

\[\begin{aligned}\gamma_{234}=\gamma_{27} \gamma_{57}\\ \gamma_{235}=\gamma_{27} \gamma_{29}\\ \gamma_{236}=\varsigma  \gamma_{214}\end{aligned}\]

\[\begin{aligned}\gamma_{237}=\gamma_{0} \gamma_{114}\\ \gamma_{238}=- \gamma_{50}\\ \gamma_{239}=\gamma_{238} + \gamma_{31}\end{aligned}\]

\[\begin{aligned}\gamma_{240}=- \gamma_{25} - \gamma_{61}\\ \gamma_{241}=4 \gamma_{13} w\\ \gamma_{242}=- \gamma_{112}\end{aligned}\]

\[\begin{aligned}\gamma_{243}=- 4 \gamma_{236}\\ \gamma_{244}=\gamma_{108} \gamma_{158}\\ \gamma_{245}=\gamma_{17} w\end{aligned}\]

\[\begin{aligned}\gamma_{246}=\gamma_{245} \gamma_{6}\\ \gamma_{247}=\gamma_{108} + \gamma_{129}\\ \gamma_{248}=\gamma_{43} \gamma_{45}\end{aligned}\]

\[\begin{aligned}\gamma_{249}=\gamma_{248} + \gamma_{97}\\ \gamma_{250}=\gamma_{108} \gamma_{163}\\ \gamma_{251}=\gamma_{25} + \gamma_{250} + \gamma_{89}\end{aligned}\]

\[\begin{aligned}\gamma_{252}=- \gamma_{230} + \gamma_{26} + \gamma_{34} \gamma_{43}\\ \gamma_{253}=- \gamma_{209} + \gamma_{30}\\ \gamma_{254}=- \gamma_{215} + \gamma_{253}\end{aligned}\]

\[\begin{aligned}\gamma_{255}=7 \gamma_{13}\\ \gamma_{256}=- 18 \zeta  \gamma_{12}\\ \gamma_{257}=- \gamma_{13} \gamma_{21} + \gamma_{21}\end{aligned}\]

\[\begin{aligned}\gamma_{258}=\begin{aligned}& 18 \eta  \zeta - 6 \eta + 6 \zeta - \gamma_{124}- 18 \gamma_{14}- \gamma_{189}- \gamma_{207}- \gamma_{23}- \gamma_{255}\\&- \gamma_{256}- \gamma_{257}\end{aligned}\end{aligned}\]

\[\begin{aligned}\gamma_{259}=\gamma_{258} \gamma_{8}\\ \gamma_{260}=\gamma_{259} \gamma_{5}\\ \gamma_{261}=- \gamma_{29} \gamma_{92}\end{aligned}\]

\[\begin{aligned}\gamma_{262}=13 \gamma_{0}\\ \gamma_{263}=6 \gamma_{12}\\ \gamma_{264}=6 \gamma_{13}\end{aligned}\]

\[\begin{aligned}\gamma_{265}=- 4 \gamma_{70}\\ \gamma_{266}=2 \gamma_{47}\\ \gamma_{267}=34 \gamma_{14}\end{aligned}\]

\[\begin{aligned}\gamma_{268}=\varsigma  \gamma_{10}\\ \gamma_{269}=\varsigma  \gamma_{50}\\ \gamma_{270}=56 \gamma_{269} + \gamma_{79}\end{aligned}\]

\[\begin{aligned}\gamma_{271}=- \gamma_{46}\\ \gamma_{272}=4 \gamma_{164}\\ \gamma_{273}=\gamma_{271} + \gamma_{272}\end{aligned}\]

\[\begin{aligned}\gamma_{274}=6 \gamma_{15}\\ \gamma_{275}=- \gamma_{274}\\ \gamma_{276}=3 \gamma_{218}\end{aligned}\]

\[\begin{aligned}\gamma_{277}=\gamma_{275} + \gamma_{276} - 17 \gamma_{66} - 47 \gamma_{71}\\ \gamma_{278}=- \gamma_{197}\\ \gamma_{279}=\varsigma  \gamma_{263}\end{aligned}\]

\[\begin{aligned}\gamma_{280}=6 \gamma_{29}\\ \gamma_{281}=3 \gamma_{108}\\ \gamma_{282}=6 \gamma_{10}\end{aligned}\]

\[\begin{aligned}\gamma_{283}=18 \gamma_{50}\\ \gamma_{284}=15 \gamma_{71}\\ \gamma_{285}=3 \gamma_{107}\end{aligned}\]

\[\begin{aligned}\gamma_{286}=- \gamma_{285}\\ \gamma_{287}=\gamma_{286} + \gamma_{81}\\ \gamma_{288}=- 3 \gamma_{21} \gamma_{31}\end{aligned}\]

\[\begin{aligned}\gamma_{289}=12 \gamma_{0}\\ \gamma_{290}=10 \gamma_{50}\\ \gamma_{291}=3 \gamma_{164}\end{aligned}\]

\[\begin{aligned}\gamma_{292}=13 \gamma_{108}\\ \gamma_{293}=6 \gamma_{114}\\ \gamma_{294}=- 6 \eta  \gamma_{13} - 6 \zeta  \gamma_{100} \gamma_{12} + \gamma_{293}\end{aligned}\]

\[\begin{aligned}\gamma_{295}=- \gamma_{248}\\ \gamma_{296}=\gamma_{295} + \gamma_{97}\\ \gamma_{297}=- \gamma_{56}\end{aligned}\]

\[\begin{aligned}\gamma_{298}=\gamma_{131} + \gamma_{297}\\ \gamma_{299}=- 39 \zeta  \gamma_{100} \gamma_{35} - 3 \gamma_{0} \gamma_{100} \gamma_{21}\\ \gamma_{300}=- \gamma_{82}\end{aligned}\]

\[\begin{aligned}\gamma_{301}=- \zeta  \gamma_{35}\\ \gamma_{302}=- 6 \varsigma  \gamma_{0} \gamma_{35} + \gamma_{223} + \gamma_{300} + \gamma_{301} + \gamma_{48} + \gamma_{86}\end{aligned}\]

\[\begin{aligned}\gamma_{303}=\gamma_{114} \gamma_{55} - \gamma_{149} + \gamma_{224} + \gamma_{284}\\ \gamma_{304}=12 \zeta \\ \gamma_{305}=13 \gamma_{114}\end{aligned}\]

\[\begin{aligned}\gamma_{306}=27 \gamma_{10}\\ \gamma_{307}=\gamma_{130} + \gamma_{154}\\ \gamma_{308}=\gamma_{169} + \gamma_{195}\end{aligned}\]

\[\begin{aligned}\gamma_{309}=3 \gamma_{195}\\ \gamma_{310}=- \gamma_{309}\\ \gamma_{311}=3 \gamma_{71}\end{aligned}\]

\[\begin{aligned}\gamma_{312}=\gamma_{114} \gamma_{92}\\ \gamma_{313}=6 \gamma_{0}\\ \gamma_{314}=34 \gamma_{1}\end{aligned}\]

\[\begin{aligned}\gamma_{315}=4 \gamma_{23}\\ \gamma_{316}=- 3 \gamma_{36}\\ \gamma_{317}=59 \gamma_{232} + \gamma_{316}\end{aligned}\]

\[\begin{aligned}\gamma_{318}=6 \gamma_{50}\\ \gamma_{319}=\gamma_{100} \gamma_{193}\\ \gamma_{320}=4 \gamma_{234}\end{aligned}\]

\[\begin{aligned}\gamma_{321}=\gamma_{149} + \gamma_{192} \gamma_{55} + 12 \gamma_{220} + \gamma_{222} + \gamma_{318} + \gamma_{319} - \gamma_{320}\\ \gamma_{322}=\gamma_{23} \gamma_{43}\end{aligned}\]

\[\begin{aligned}\gamma_{323}=28 \gamma_{164}\\ \gamma_{324}=\gamma_{100} \gamma_{323} + \gamma_{162}\\ \gamma_{325}=- \gamma_{201}\end{aligned}\]

\[\begin{aligned}\gamma_{326}=\gamma_{101} + \gamma_{108} \gamma_{255} + \gamma_{325}\\ \gamma_{327}=- 6 \zeta  \gamma_{35} + \gamma_{0}\\ \gamma_{328}=- \gamma_{13}\end{aligned}\]

\[\begin{aligned}\gamma_{329}=\gamma_{12} + \gamma_{328}\\ \gamma_{330}=- 18 \zeta  \gamma_{21} + \gamma_{329} + 6 \gamma_{35} + \gamma_{84}\end{aligned}\]

\[\begin{aligned}\gamma_{331}=- \gamma_{182} + \gamma_{23} - \gamma_{256} - \gamma_{283} - \gamma_{327} - \gamma_{330} - \gamma_{94}\\ \gamma_{332}=\gamma_{331} \gamma_{9}\end{aligned}\]

\[\begin{aligned}\gamma_{333}=\gamma_{2} \gamma_{332} \gamma_{5}\\ \gamma_{334}=- 12 \zeta  \gamma_{35}\\ \gamma_{335}=\gamma_{264} \gamma_{35}\end{aligned}\]

\[\begin{aligned}\gamma_{336}=8 \gamma_{117}\\ \gamma_{337}=\gamma_{10} \gamma_{43}\\ \gamma_{338}=- \gamma_{66}\end{aligned}\]

\[\begin{aligned}\gamma_{339}=- \gamma_{195}\\ \gamma_{340}=14 \gamma_{23}\\ \gamma_{341}=4 \gamma_{117}\end{aligned}\]

\[\begin{aligned}\gamma_{342}=2 \gamma_{194}\\ \gamma_{343}=\gamma_{193} + \gamma_{31}\\ \gamma_{344}=\gamma_{22} - \gamma_{304} \gamma_{67}\end{aligned}\]

\[\begin{aligned}\gamma_{345}=\gamma_{10} + \gamma_{205} - \gamma_{77}\\ \gamma_{346}=- 6 \gamma_{48}\\ \gamma_{347}=13 \gamma_{101}\end{aligned}\]

\[\begin{aligned}\gamma_{348}=6 \gamma_{131}\\ \gamma_{349}=3 \gamma_{50}\\ \gamma_{350}=\gamma_{1} + \gamma_{201} + \gamma_{271}\end{aligned}\]

\[\begin{aligned}\gamma_{351}=\gamma_{190} + \gamma_{194} + \gamma_{312}\\ \gamma_{352}=- \gamma_{114}\\ \gamma_{353}=\gamma_{14} + \gamma_{352}\end{aligned}\]

\[\begin{aligned}\gamma_{354}=\gamma_{15} - 12 \gamma_{65}\\ \gamma_{355}=3 \gamma_{105}\\ \gamma_{356}=- \gamma_{355}\end{aligned}\]

\[\begin{aligned}\gamma_{357}=- \gamma_{171} - \gamma_{313} \gamma_{57} + \gamma_{356} + \gamma_{53} + \gamma_{55} + \gamma_{96}\\ \gamma_{358}=10 \gamma_{14}\\ \gamma_{359}=2 \gamma_{57}\end{aligned}\]

\[\begin{aligned}\gamma_{360}=15 \gamma_{228}\\ \gamma_{361}=25 \gamma_{74}\\ \gamma_{362}=\gamma_{10} \gamma_{100}\end{aligned}\]

\[\begin{aligned}\gamma_{363}=8 \gamma_{362}\\ \gamma_{364}=- 6 \gamma_{71}\\ \gamma_{365}=3 \gamma_{194}\end{aligned}\]

\[\begin{aligned}\gamma_{366}=2 \eta \\ \gamma_{367}=\gamma_{100} \gamma_{366}\\ \gamma_{368}=3 \gamma_{192} + \gamma_{367}\end{aligned}\]

\[\begin{aligned}\gamma_{369}=5 \gamma_{1}\\ \gamma_{370}=\gamma_{280} - \gamma_{369}\\ \gamma_{371}=\gamma_{100} \gamma_{14}\end{aligned}\]

\[\begin{aligned}\gamma_{372}=44 \gamma_{68}\\ \gamma_{373}=22 \gamma_{107}\\ \gamma_{374}=11 \gamma_{71}\end{aligned}\]

\[\begin{aligned}\gamma_{375}=\gamma_{188} \gamma_{57}\\ \gamma_{376}=22 \gamma_{58}\\ \gamma_{377}=110 \gamma_{111}\end{aligned}\]

\[\begin{aligned}\gamma_{378}=\gamma_{2} + \gamma_{81}\\ \gamma_{379}=- 6 \zeta  \gamma_{100} \gamma_{42} + 12 \gamma_{216}\\ \gamma_{380}=14 \gamma_{50}\end{aligned}\]

\[\begin{aligned}\gamma_{381}=14 \gamma_{195}\\ \gamma_{382}=6 \gamma_{179}\\ \gamma_{383}=7 \gamma_{105}\end{aligned}\]

\[\begin{aligned}\gamma_{384}=68 \gamma_{111}\\ \gamma_{385}=\gamma_{0} + \gamma_{114}\\ \gamma_{386}=\eta  + 1\end{aligned}\]

\[\begin{aligned}\gamma_{387}=17 \zeta  \gamma_{12} + \gamma_{148} - \gamma_{257} - \gamma_{290} - \gamma_{329} - \gamma_{343} - \gamma_{358}\\ \gamma_{388}=\gamma_{386} \gamma_{387}\end{aligned}\]

\[\begin{aligned}\gamma_{389}=\gamma_{2} \gamma_{388} \gamma_{7}\\ \gamma_{390}=- 4 \gamma_{167}\\ \gamma_{391}=\gamma_{81} + \gamma_{95}\end{aligned}\]

\[\begin{aligned}\gamma_{392}=\gamma_{13} - \gamma_{337}\\ \gamma_{393}=- \gamma_{160} + \gamma_{214}\\ \gamma_{394}=- 5 \gamma_{29}\end{aligned}\]

\[\begin{aligned}\gamma_{395}=\gamma_{114} + \gamma_{74}\\ \gamma_{396}=- \gamma_{198}\\ \gamma_{397}=\gamma_{396} + \gamma_{77}\end{aligned}\]

\[\begin{aligned}\gamma_{398}=\gamma_{132} - \gamma_{323}\\ \gamma_{399}=3 \gamma_{222}\\ \gamma_{400}=\gamma_{100} \gamma_{13}\end{aligned}\]

\[\begin{aligned}\gamma_{401}=\gamma_{400} \gamma_{84}\\ \gamma_{402}=12 \gamma_{63}\\ \gamma_{403}=46 \gamma_{221}\end{aligned}\]

\[\begin{aligned}\gamma_{404}=\gamma_{194} + \gamma_{353}\\ \gamma_{405}=- 2 \gamma_{100} \gamma_{35} + \gamma_{146} + \gamma_{212} + \gamma_{70} + 24 \gamma_{73}\end{aligned}\]

\[\begin{aligned}\gamma_{406}=\gamma_{154} + \gamma_{159}\\ \gamma_{407}=\gamma_{185} + \gamma_{66}\\ \gamma_{408}=\gamma_{124} + \gamma_{31}\end{aligned}\]

\[\begin{aligned}\gamma_{409}=\gamma_{27} - \gamma_{318}\\ \gamma_{410}=- \gamma_{57}\\ \gamma_{411}=\gamma_{4} \gamma_{5}\end{aligned}\]

\[\begin{aligned}\gamma_{412}=\varsigma ^{3}\\ \gamma_{413}=\gamma_{21} \gamma_{412}\\ \gamma_{414}=- \gamma_{413}\end{aligned}\]

\[\begin{aligned}\gamma_{415}=- \gamma_{12} \gamma_{412}\\ \gamma_{416}=2 \gamma_{15}\\ \gamma_{417}=\gamma_{35} \gamma_{412}\end{aligned}\]

\[\begin{aligned}\gamma_{418}=\gamma_{192} + \gamma_{309} + \gamma_{391} + \gamma_{417}\\ \gamma_{419}=\gamma_{130} + \gamma_{33} \gamma_{412}\\ \gamma_{420}=\gamma_{412} \gamma_{70}\end{aligned}\]

\[\begin{aligned}\gamma_{421}=\gamma_{286} + \gamma_{413} + 2 \gamma_{420}\\ \gamma_{422}=\varsigma  \gamma_{349} + \gamma_{356} + \gamma_{385}\end{aligned}\]

\[\begin{aligned}\gamma_{423}=- \gamma_{100} \gamma_{95} + \gamma_{309} + \gamma_{96}\\ \gamma_{424}=2 \gamma_{50}\end{aligned}\]

\endgroup
The twenty entries are listed by matrix row. Each expression retains
the auxiliary quantities defined above.

\subsection{Matrix entries}
\begingroup\small
\setlength{\abovedisplayskip}{8pt}
\setlength{\belowdisplayskip}{8pt}
\RenewDocumentEnvironment{aligned}{O{t}+b}{%
  \begingroup
  \let\aligned\suppaligned
  \let\endaligned\endsuppaligned
  \suppaligned[#1]#2\endsuppaligned
  \endgroup
}{}

\subsubsection*{Row 1}
\[\mathsf A_{1,1}=\begin{aligned}&- \gamma_{20} w\\&- \gamma_{92} \left(\begin{aligned}& 55 \varsigma  \gamma_{13} \gamma_{21}+ 21 \varsigma  \gamma_{13} \gamma_{33}+ 55 \varsigma  \gamma_{13} \gamma_{35}+ 2 \varsigma  \gamma_{21}\\&+ 18 \varsigma  \gamma_{33}+ 18 \varsigma  \gamma_{35}- 10 \gamma_{1}+ 16 \gamma_{12} \gamma_{13}+ 9 \gamma_{12} \gamma_{31}\\&+ 33 \gamma_{13} \gamma_{21}+ 3 \gamma_{13} \gamma_{33}+ 15 \gamma_{13} \gamma_{35}- \gamma_{16}- \gamma_{21}- \gamma_{22}\\&- \gamma_{24}- \gamma_{26}- \gamma_{30}- \gamma_{32}+ \gamma_{33}- \gamma_{34}- \gamma_{37}+ \gamma_{39}+ \gamma_{41}+ \gamma_{44}\\&- 3 \gamma_{45}- 7 \gamma_{46}- 11 \gamma_{47}- 11 \gamma_{48}- 21 \gamma_{49}- \gamma_{52}+ \gamma_{53}- \gamma_{54}\\&- \gamma_{56}- \gamma_{57} \gamma_{64}- \gamma_{59}- \gamma_{61}- 4 \gamma_{63}- 8 \gamma_{65}- 11 \gamma_{66}\\&- 11 \gamma_{68}- 18 \gamma_{69}- 31 \gamma_{71}- 31 \gamma_{72}- 66 \gamma_{73}- \gamma_{80}- \gamma_{87}\\&- \gamma_{91}\end{aligned}\right)\end{aligned}\]

\[\mathsf A_{1,2}=\begin{aligned}& 8 \gamma_{0} \gamma_{11} \gamma_{12} \gamma_{18} \gamma_{5} \gamma_{6} \gamma_{8} \gamma_{9} \gamma_{93}\\&- \gamma_{142} \left(\begin{aligned}& \eta  \varsigma  \gamma_{31}+ 35 \varsigma  \zeta  \gamma_{21}+ 4 \varsigma  \zeta  \gamma_{33}+ 37 \varsigma  \zeta  \gamma_{35}+ \varsigma  \gamma_{0} \gamma_{35}\\&+ \varsigma  \gamma_{21} \gamma_{31}+ 20 \zeta  \gamma_{100} \gamma_{33}+ 14 \zeta  \gamma_{100} \gamma_{35}+ 2 \zeta  \gamma_{100} \gamma_{42}\\&- 10 \gamma_{101}- \gamma_{102}- \gamma_{103}- \gamma_{104}- \gamma_{106}- 4 \gamma_{107}\\&- \gamma_{109} \gamma_{13}- \gamma_{110}- \gamma_{112}- \gamma_{116}- \gamma_{119}+ 23 \gamma_{12} \gamma_{31}\\&- \gamma_{123}- \gamma_{128}- \gamma_{133}- \gamma_{136}- \gamma_{139}- 8 \gamma_{66}- \gamma_{94}- \gamma_{95}\\&- \gamma_{96}- \gamma_{97}- \gamma_{99}\end{aligned}\right)\\&- \gamma_{183} \left(\begin{aligned}& 6 \eta  \varsigma  \gamma_{31}+ 3 \eta  \gamma_{144}+ 3 \eta  \gamma_{31}+ \varsigma  \gamma_{12} \gamma_{27}+ 47 \varsigma  \gamma_{13} \gamma_{21}\\&+ 30 \varsigma  \gamma_{13} \gamma_{33}+ 91 \varsigma  \gamma_{13} \gamma_{35}+ 2 \varsigma  \gamma_{21} \gamma_{25}+ 3 \varsigma  \gamma_{21} \gamma_{27}\\&+ 5 \varsigma  \gamma_{21} \gamma_{31}+ 2 \varsigma  \gamma_{25} \gamma_{35}+ 4 \varsigma  \gamma_{31} \gamma_{33}+ 11 \varsigma  \gamma_{31} \gamma_{35}\\&- \zeta  \gamma_{162}+ 3 \gamma_{0} \gamma_{100} \gamma_{12}+ \gamma_{0} \gamma_{100} \gamma_{145}+ 7 \gamma_{100} \gamma_{12} \gamma_{31}\\&+ 34 \gamma_{100} \gamma_{13} \gamma_{33}+ 9 \gamma_{100} \gamma_{13} \gamma_{35}+ 16 \gamma_{100} \gamma_{13} \gamma_{42}\\&+ 11 \gamma_{100} \gamma_{21} \gamma_{31}- \gamma_{100} \gamma_{27} \gamma_{33}+ 11 \gamma_{100} \gamma_{31} \gamma_{33}\\&+ 15 \gamma_{100} \gamma_{31} \gamma_{35}- \gamma_{101}- \gamma_{104}- \gamma_{108}+ 25 \gamma_{12} \gamma_{27}\\&+ 25 \gamma_{12} \gamma_{31}+ 10 \gamma_{13} \gamma_{21}+ 4 \gamma_{13} \gamma_{35}- \gamma_{131} \gamma_{163}- \gamma_{143}\\&- \gamma_{144} \gamma_{21}+ 2 \gamma_{144}+ \gamma_{146}+ \gamma_{147}- \gamma_{148}- \gamma_{150}\\&- 10 \gamma_{151}- \gamma_{152}- \gamma_{153}+ \gamma_{156}- \gamma_{157}- \gamma_{160}- 3 \gamma_{161}\\&- 9 \gamma_{165}- 10 \gamma_{166}- 14 \gamma_{167}- 33 \gamma_{168}- \gamma_{172}- \gamma_{174}\\&- \gamma_{176}- \gamma_{181}+ 40 \gamma_{21} \gamma_{27}- 25 \gamma_{23}- \gamma_{26} \gamma_{43}- 25 \gamma_{26}\\&+ \gamma_{27} \gamma_{33}+ 4 \gamma_{27} \gamma_{35}- \gamma_{33} \gamma_{55}- 4 \gamma_{38}- \gamma_{45}- 4 \gamma_{48}\\&- 40 \gamma_{49}- \gamma_{54}- \gamma_{55} \gamma_{60}- \gamma_{57} \gamma_{82}- 22 \gamma_{65}- 5 \gamma_{66}\\&- 68 \gamma_{68}- 55 \gamma_{69}+ \gamma_{70}- 20 \gamma_{73}- \gamma_{82}\end{aligned}\right)\end{aligned}\]

\[\mathsf A_{1,3}=\begin{aligned}&- \gamma_{20} \gamma_{93}\\&+ \gamma_{204} \left(\begin{aligned}& 17 \varsigma  \gamma_{15}+ 14 \gamma_{1}+ 8 \gamma_{100} \gamma_{14}- 14 \gamma_{107}- 4 \gamma_{111}- \gamma_{113}\\&+ 10 \gamma_{114}+ \gamma_{115}+ \gamma_{134}- \gamma_{138}- 20 \gamma_{14}+ 8 \gamma_{15}- 15 \gamma_{154}\\&+ \gamma_{167}+ \gamma_{17}+ \gamma_{185}+ \gamma_{186}- \gamma_{187}+ \gamma_{188}- \gamma_{189}- \gamma_{190}\\&+ \gamma_{191}+ 2 \gamma_{192}+ 4 \gamma_{193}- 20 \gamma_{194}- 19 \gamma_{195}- \gamma_{196}\\&+ \gamma_{197}+ \gamma_{198}+ \gamma_{199}+ \gamma_{200}+ \gamma_{203}- 15 \gamma_{29}+ \gamma_{62}- 4 \gamma_{66}\\&- \gamma_{67}+ 19 \gamma_{71}\end{aligned}\right)\\&+ \gamma_{241} \left(\begin{aligned}& 2 \eta  \gamma_{144}+ 12 \varsigma  \gamma_{213}- 8 \varsigma  \gamma_{26}- \zeta  \gamma_{225}- 3 \gamma_{0} \gamma_{131}\\&+ \gamma_{100} \gamma_{46}- 7 \gamma_{101} \gamma_{31}- \gamma_{101}+ \gamma_{105}+ \gamma_{107}+ \gamma_{12} \gamma_{144}\\&+ \gamma_{131}+ 55 \gamma_{138}+ \gamma_{144}- \gamma_{146} \gamma_{158}+ \gamma_{146}+ 7 \gamma_{149}\\&- 3 \gamma_{151}+ 3 \gamma_{153}- 4 \gamma_{155}+ \gamma_{161}- 16 \gamma_{165}+ 7 \gamma_{166}\\&- 55 \gamma_{167}+ \gamma_{181}+ \gamma_{188} \gamma_{77}+ 10 \gamma_{193}+ \gamma_{205}+ \gamma_{206}\\&+ \gamma_{209}+ \gamma_{211}+ \gamma_{212}+ \gamma_{213}+ 3 \gamma_{214}+ \gamma_{215}+ 11 \gamma_{216}\\&+ 21 \gamma_{217}+ 34 \gamma_{218}- \gamma_{219}- 37 \gamma_{220}- 34 \gamma_{221}- 15 \gamma_{222}\\&- 11 \gamma_{223}- 11 \gamma_{224}- 9 \gamma_{226}- \gamma_{227} \gamma_{29}- \gamma_{227} \gamma_{57}\\&+ 8 \gamma_{229}- 16 \gamma_{23}+ 9 \gamma_{230}+ \gamma_{231} \gamma_{29}+ 10 \gamma_{232}\\&+ \gamma_{233} \gamma_{27}+ 22 \gamma_{234}+ 23 \gamma_{235}+ 33 \gamma_{236}+ 33 \gamma_{237}\\&+ \gamma_{239}+ \gamma_{240}- 9 \gamma_{26}- 7 \gamma_{27} \gamma_{35}- 3 \gamma_{34}+ \gamma_{37}+ 11 \gamma_{38}\\&- 18 \gamma_{46}- 15 \gamma_{48}- 33 \gamma_{49}- 10 \gamma_{58}- 23 \gamma_{65}- 37 \gamma_{68}+ \gamma_{70}\\&- 22 \gamma_{72}+ \gamma_{76}+ 33 \gamma_{88}+ \gamma_{94}\end{aligned}\right)\end{aligned}\]

\[\mathsf A_{1,4}=\begin{aligned}& 18 \gamma_{0} \gamma_{67}+ \gamma_{121}- \gamma_{126}+ \gamma_{131} \gamma_{158}- 32 \gamma_{138}+ 23 \gamma_{165}\\&+ 32 \gamma_{167}+ \gamma_{170}+ \gamma_{174}+ 2 \gamma_{179}+ \gamma_{19} \gamma_{246}+ \gamma_{210}\\&+ 18 \gamma_{220}- 23 \gamma_{226}+ 23 \gamma_{23}+ \gamma_{238}+ \gamma_{242}+ \gamma_{243}- \gamma_{244}\\&+ \gamma_{247}+ \gamma_{249}+ \gamma_{251}+ \gamma_{252}+ \gamma_{254}- 23 \gamma_{36}+ 14 \gamma_{58}\\&- 28 \gamma_{66}+ 14 \gamma_{73}+ \gamma_{87}- \gamma_{94}\end{aligned}\]

\subsubsection*{Row 2}
\[\mathsf A_{2,1}=\begin{aligned}& \gamma_{2} \gamma_{260} w\\&+ \gamma_{92} \left(\begin{aligned}&- 59 \gamma_{1}+ 21 \gamma_{111}+ 27 \gamma_{117}- \gamma_{122}- 8 \gamma_{138}+ \gamma_{187}\\&+ 26 \gamma_{193}- 11 \gamma_{195}+ \gamma_{201}- \gamma_{216}+ \gamma_{225}- 39 \gamma_{23}+ \gamma_{249}\\&+ \gamma_{261}- \gamma_{262}- \gamma_{263}+ \gamma_{264}+ \gamma_{265}+ \gamma_{266}+ \gamma_{267}\\&- 12 \gamma_{268}+ \gamma_{270}+ \gamma_{273}+ \gamma_{277}+ 24 \gamma_{29}- \gamma_{35}+ 48 \gamma_{50}\\&+ 4 \gamma_{73}- 59 \gamma_{74}+ 12 \gamma_{77}- \gamma_{84}\end{aligned}\right)\end{aligned}\]

\[\mathsf A_{2,2}=\begin{aligned}& \gamma_{142} \left(\begin{aligned}& \varsigma  \gamma_{282}- \varsigma  \gamma_{283}- \zeta  \gamma_{134}+ 29 \gamma_{1}+ 2 \gamma_{101}+ \gamma_{111}+ \gamma_{114}\\&+ 21 \gamma_{15}+ \gamma_{163}- \gamma_{193}+ 5 \gamma_{194}+ \gamma_{195}- \gamma_{200}+ \gamma_{208}\\&- \gamma_{264}- \gamma_{267}+ \gamma_{278}- \gamma_{279}- \gamma_{280}+ \gamma_{281}+ \gamma_{282}+ \gamma_{284}\\&+ \gamma_{287}- 28 \gamma_{50}+ \gamma_{66}+ 2 \gamma_{67}+ 23 \gamma_{74}\end{aligned}\right)\\&+ \gamma_{183} \left(\begin{aligned}& 12 \eta  \varsigma  \gamma_{13}+ 26 \eta  \varsigma  \gamma_{31}+ 98 \eta  \gamma_{31}+ 5 \varsigma  \zeta  \gamma_{35}+ 59 \varsigma  \gamma_{12} \gamma_{13}\\&+ 36 \varsigma  \gamma_{12} \gamma_{31}+ 32 \varsigma  \gamma_{13} \gamma_{21}+ 23 \zeta  \gamma_{100} \gamma_{21}\\&+ 22 \zeta  \gamma_{100} \gamma_{33}+ 10 \gamma_{0} \gamma_{100} \gamma_{35}- 57 \gamma_{1}+ 68 \gamma_{12} \gamma_{31}\\&- \gamma_{135}- \gamma_{137}- \gamma_{139}- 24 \gamma_{165}- 36 \gamma_{195}- \gamma_{199}- \gamma_{209}\\&- \gamma_{216}- 29 \gamma_{221}- 72 \gamma_{23}- \gamma_{252}- 16 \gamma_{27}- \gamma_{275}- \gamma_{288}\\&- \gamma_{289}+ \gamma_{290}- \gamma_{291}- \gamma_{292}- \gamma_{294}- \gamma_{296}- \gamma_{298}- \gamma_{299}\\&- \gamma_{30}- \gamma_{302}- \gamma_{303}+ \gamma_{31} \gamma_{35}+ 26 \gamma_{31}- 11 \gamma_{36}- 48 \gamma_{46}\\&- 71 \gamma_{66}- 21 \gamma_{69}- 12 \gamma_{74}\end{aligned}\right)\\&- 8 \gamma_{23} \gamma_{260} \gamma_{93}\end{aligned}\]

\[\mathsf A_{2,3}=\begin{aligned}& 8 \eta  \gamma_{0} \gamma_{258} \gamma_{5} \gamma_{8} \gamma_{93}\\&- \gamma_{204} \left(\begin{aligned}& 12 \eta  \gamma_{100}- \varsigma  \gamma_{304}- \gamma_{100} \gamma_{306}+ 11 \gamma_{117}- 25 \gamma_{13}\\&- 25 \gamma_{14}+ 25 \gamma_{192}- 25 \gamma_{194}+ 13 \gamma_{228}+ \gamma_{262}+ \gamma_{29}\\&+ \gamma_{304}+ \gamma_{305}+ \gamma_{306}+ \gamma_{307}+ \gamma_{308}- 12 \gamma_{57}- \gamma_{71}\\&+ 11 \gamma_{74}- 13 \gamma_{77}\end{aligned}\right)\\&- \gamma_{241} \left(\begin{aligned}& 6 \gamma_{100} \gamma_{14}+ \gamma_{100} \gamma_{315}- 6 \gamma_{105}+ 34 \gamma_{107}- 9 \gamma_{111}\\&- \gamma_{113} \gamma_{35}+ \gamma_{125}- \gamma_{137}- 21 \gamma_{138}- 48 \gamma_{165}+ \gamma_{166}\\&+ 21 \gamma_{167}+ \gamma_{173}+ 26 \gamma_{179}+ 59 \gamma_{193}+ 39 \gamma_{218}- 7 \gamma_{219}\\&- 39 \gamma_{221}- 3 \gamma_{226}- 48 \gamma_{23}+ \gamma_{243}- \gamma_{250}- 8 \gamma_{27}+ \gamma_{272}\\&+ \gamma_{274}+ \gamma_{289} \gamma_{67}- \gamma_{292}- \gamma_{293}+ 13 \gamma_{31}+ \gamma_{310}+ \gamma_{311}\\&+ \gamma_{312}- \gamma_{313}- \gamma_{314}+ \gamma_{317}+ \gamma_{321}- 26 \gamma_{46}+ \gamma_{47}+ \gamma_{48}\\&- 4 \gamma_{49}+ \gamma_{58}- 9 \gamma_{66}+ 4 \gamma_{72}+ \gamma_{73}+ \gamma_{91}\end{aligned}\right)\end{aligned}\]

\[\mathsf A_{2,4}=\begin{aligned}& 29 \eta  \gamma_{31}+ 34 \zeta  \gamma_{100} \gamma_{21}+ \zeta  \gamma_{100} \gamma_{33}+ 29 \zeta  \gamma_{100} \gamma_{35}\\&- 21 \gamma_{105}- \gamma_{113} \gamma_{29}- \gamma_{113} \gamma_{67}- \gamma_{118}+ 21 \gamma_{12} \gamma_{13}\\&+ 7 \gamma_{12} \gamma_{31}- 28 \gamma_{23}- \gamma_{245} \gamma_{259}- \gamma_{278}- \gamma_{294}+ 7 \gamma_{31}\\&- \gamma_{313}- \gamma_{314}- \gamma_{322}- \gamma_{324}- \gamma_{326}- \gamma_{46}\end{aligned}\]

\subsubsection*{Row 3}
\[\mathsf A_{3,1}=\begin{aligned}&- \gamma_{333} w\\&- \gamma_{92} \left(\begin{aligned}& 8 \eta  \gamma_{13}+ 9 \eta  \gamma_{31}+ 59 \varsigma  \zeta  \gamma_{35}+ 14 \gamma_{0} \gamma_{21}- 18 \gamma_{1}+ \gamma_{110}\\&- 56 \gamma_{111}- 12 \gamma_{127}- 27 \gamma_{138}- 25 \gamma_{164}+ \gamma_{175}- \gamma_{203}\\&- 3 \gamma_{23}- \gamma_{239}- \gamma_{271}- \gamma_{277}- 10 \gamma_{29}- \gamma_{330}- \gamma_{334}\\&- \gamma_{335}+ \gamma_{336}- \gamma_{337}- \gamma_{51}- 24 \gamma_{67}- \gamma_{80}- \gamma_{99}\end{aligned}\right)\end{aligned}\]

\[\mathsf A_{3,2}=\begin{aligned}& 8 \gamma_{0} \gamma_{12} \gamma_{331} \gamma_{5} \gamma_{9} \gamma_{93}\\&- \gamma_{142} \left(\begin{aligned}& 12 \gamma_{101}- 25 \gamma_{107}+ 25 \gamma_{108}+ 41 \gamma_{111}+ 14 \gamma_{15}- 27 \gamma_{232}\\&+ \gamma_{233}+ \gamma_{268}+ 39 \gamma_{269}+ 13 \gamma_{29}+ \gamma_{305}+ \gamma_{328}+ \gamma_{338}\\&+ \gamma_{339}- \gamma_{340}- \gamma_{341}+ \gamma_{342}+ \gamma_{343}+ \gamma_{344}+ \gamma_{345}\\&- 50 \gamma_{71}- 37 \gamma_{74}\end{aligned}\right)\\&- \gamma_{183} \left(\begin{aligned}& 6 \zeta  \gamma_{35}+ \gamma_{0} \gamma_{102}- \gamma_{101} \gamma_{264}+ 26 \gamma_{107}+ \gamma_{108} \gamma_{92}\\&- 109 \gamma_{111}+ \gamma_{131} \gamma_{264}- \gamma_{134}- 68 \gamma_{138}- 39 \gamma_{165}\\&- 29 \gamma_{167}+ \gamma_{176}+ 34 \gamma_{179}+ \gamma_{184}+ \gamma_{211}+ 4 \gamma_{217}\\&- 9 \gamma_{220}- 48 \gamma_{221}+ \gamma_{223}+ \gamma_{254}+ \gamma_{26}- 31 \gamma_{269}- \gamma_{272}\\&+ \gamma_{300}+ \gamma_{303}+ \gamma_{317}- \gamma_{335}+ \gamma_{346}- \gamma_{347}- \gamma_{348}- \gamma_{349}\\&+ \gamma_{350}+ \gamma_{351}+ \gamma_{353}+ \gamma_{354}+ \gamma_{357}+ 6 \gamma_{38}+ 14 \gamma_{63}\\&+ 78 \gamma_{66}+ 105 \gamma_{69}+ 23 \gamma_{73}+ \gamma_{90}\end{aligned}\right)\end{aligned}\]

\[\mathsf A_{3,3}=\begin{aligned}& \gamma_{204} \left(\begin{aligned}& 15 \gamma_{10}+ 5 \gamma_{107}+ \gamma_{114}+ 24 \gamma_{117}- \gamma_{154}+ \gamma_{158}- \gamma_{159}\\&- \gamma_{231}- 21 \gamma_{268}+ 7 \gamma_{269}+ \gamma_{279}+ \gamma_{339}- \gamma_{358}- \gamma_{359}\\&+ \gamma_{360}- \gamma_{361}- \gamma_{363}+ \gamma_{364}- \gamma_{365}+ \gamma_{368}+ \gamma_{370}+ \gamma_{64}\end{aligned}\right)\\&+ \gamma_{241} \left(\begin{aligned}& 24 \varsigma  \zeta  \gamma_{33}+ 41 \varsigma  \zeta  \gamma_{35}+ 79 \varsigma  \gamma_{0} \gamma_{12}+ 119 \varsigma  \gamma_{0} \gamma_{21}\\&+ 23 \zeta  \gamma_{100} \gamma_{33}+ 10 \gamma_{0} \gamma_{100} \gamma_{12}+ 25 \gamma_{0} \gamma_{21}- 8 \gamma_{105}\\&- \gamma_{114}- \gamma_{136}- 105 \gamma_{138}- 18 \gamma_{164}- 29 \gamma_{165}- 42 \gamma_{167}\\&- 14 \gamma_{217}- 24 \gamma_{221}- 4 \gamma_{226}- \gamma_{23}- 24 \gamma_{236}- \gamma_{251}\\&+ \gamma_{276}- \gamma_{299}- \gamma_{316}- \gamma_{32}- \gamma_{321}- \gamma_{327}- \gamma_{346}- \gamma_{347}\\&- \gamma_{348}+ \gamma_{371}+ \gamma_{372}+ \gamma_{373}+ \gamma_{374}+ \gamma_{375}- \gamma_{376}- \gamma_{377}\\&- \gamma_{378}- \gamma_{379}- 9 \gamma_{46}- \gamma_{51}- \gamma_{76}\end{aligned}\right)\\&- \gamma_{333} \gamma_{93}\end{aligned}\]

\[\mathsf A_{3,4}=\begin{aligned}& 6 \gamma_{101}- 29 \gamma_{107}+ \gamma_{13} \gamma_{180}+ 12 \gamma_{138}+ \gamma_{193}- \gamma_{194}\\&+ 14 \gamma_{218}- 34 \gamma_{232}+ \gamma_{245} \gamma_{332}+ \gamma_{270}+ \gamma_{324}+ \gamma_{344}\\&+ \gamma_{350}- \gamma_{380}- \gamma_{381}- \gamma_{382}+ \gamma_{383}+ \gamma_{384}+ \gamma_{385}\\&- 58 \gamma_{66}- 42 \gamma_{69}- 14 \gamma_{71}\end{aligned}\]

\subsubsection*{Row 4}
\[\mathsf A_{4,1}=\begin{aligned}&- \gamma_{389} w\\&- \gamma_{92} \left(\begin{aligned}& 20 \varsigma  \zeta  \gamma_{12}+ 20 \varsigma  \zeta  \gamma_{33}+ 6 \varsigma  \gamma_{12} \gamma_{13}+ 6 \varsigma  \gamma_{13} \gamma_{35}+ 2 \zeta  \gamma_{33}\\&+ 7 \gamma_{0} \gamma_{35}+ 5 \gamma_{12} \gamma_{13}+ 15 \gamma_{12} \gamma_{31}+ \gamma_{12}- \gamma_{124}- \gamma_{126}\\&- 6 \gamma_{127}- \gamma_{193}- \gamma_{202}+ \gamma_{21}- 18 \gamma_{216}- 8 \gamma_{236}- \gamma_{242}\\&- \gamma_{265}- \gamma_{27}- \gamma_{274}- \gamma_{288}+ \gamma_{291}- \gamma_{297}- \gamma_{325}- \gamma_{33}\\&- \gamma_{334}- \gamma_{336}- \gamma_{340}- \gamma_{35}- \gamma_{36}- \gamma_{370}- \gamma_{390}- \gamma_{391}\\&- \gamma_{392}- \gamma_{393}- \gamma_{44}- 11 \gamma_{49}- 14 \gamma_{66}- \gamma_{67} \gamma_{92}- 8 \gamma_{67}\\&- 14 \gamma_{69}- 8 \gamma_{71}- 8 \gamma_{73}- \gamma_{98}\end{aligned}\right)\end{aligned}\]

\[\mathsf A_{4,2}=\begin{aligned}& 8 \gamma_{0} \gamma_{12} \gamma_{386} \gamma_{387} \gamma_{5} \gamma_{6} \gamma_{93}\\&- \gamma_{142} \left(\begin{aligned}&- \varsigma  \gamma_{380}- \gamma_{10}+ 5 \gamma_{101}+ \gamma_{109}- 17 \gamma_{111}+ \gamma_{123}+ \gamma_{13}\\&+ \gamma_{138}- 5 \gamma_{14}- 8 \gamma_{179}- \gamma_{218}+ 29 \gamma_{23}- \gamma_{268}+ \gamma_{274}\\&+ \gamma_{296}+ \gamma_{308}+ \gamma_{322}+ \gamma_{341}- \gamma_{342}+ \gamma_{378}+ \gamma_{394}+ \gamma_{395}\\&+ \gamma_{397}+ \gamma_{398}- 34 \gamma_{50}- 4 \gamma_{67}+ \gamma_{69}+ 21 \gamma_{70}+ 17 \gamma_{71}\\&+ 13 \gamma_{73}\end{aligned}\right)\\&- \gamma_{183} \left(\begin{aligned}& 26 \eta  \gamma_{31}+ 18 \varsigma  \gamma_{0} \gamma_{33}+ \varsigma  \gamma_{12} \gamma_{13}+ 13 \varsigma  \gamma_{12} \gamma_{31}\\&+ 70 \varsigma  \gamma_{13} \gamma_{21}+ 44 \varsigma  \gamma_{13} \gamma_{35}+ 3 \varsigma  \gamma_{21} \gamma_{31}+ 2 \varsigma  \gamma_{27} \gamma_{35}\\&+ 33 \zeta  \gamma_{100} \gamma_{33}+ 6 \zeta  \gamma_{100} \gamma_{35}+ 22 \zeta  \gamma_{100} \gamma_{42}\\&+ 17 \gamma_{0} \gamma_{100} \gamma_{21}+ 17 \gamma_{0} \gamma_{100} \gamma_{33}+ 27 \gamma_{0} \gamma_{100} \gamma_{35}\\&- 13 \gamma_{1}- \gamma_{101} \gamma_{31}+ \gamma_{103}- 6 \gamma_{107}+ 82 \gamma_{12} \gamma_{31}\\&+ 3 \gamma_{13} \gamma_{35}- \gamma_{15}+ \gamma_{150}- \gamma_{151}- \gamma_{156}+ \gamma_{177}- \gamma_{201}\\&+ 50 \gamma_{21} \gamma_{31}- 51 \gamma_{226}- 73 \gamma_{23}- 14 \gamma_{230}- 23 \gamma_{236}\\&- \gamma_{253}- \gamma_{28} \gamma_{77}- \gamma_{28}- \gamma_{281} \gamma_{31}- \gamma_{302}+ \gamma_{31} \gamma_{33}- \gamma_{34}\\&- \gamma_{357}- 31 \gamma_{36}- 3 \gamma_{38}- \gamma_{393}- \gamma_{398}- \gamma_{399}- 4 \gamma_{40}\\&- \gamma_{401}- \gamma_{402}- \gamma_{403}- \gamma_{404}- \gamma_{405}- 15 \gamma_{46}- 72 \gamma_{49}\\&- 30 \gamma_{66}- 50 \gamma_{69}- 13 \gamma_{71}- 7 \gamma_{72}\end{aligned}\right)\end{aligned}\]

\[\mathsf A_{4,3}=\begin{aligned}& \gamma_{204} \left(\begin{aligned}&- \varsigma  \gamma_{290}+ 4 \gamma_{1}- 17 \gamma_{10}- 4 \gamma_{114}- 25 \gamma_{117}+ 13 \gamma_{14}- \gamma_{158}\\&+ \gamma_{16}+ \gamma_{172}- 5 \gamma_{192}+ 4 \gamma_{194}+ \gamma_{206}+ \gamma_{231}+ \gamma_{244}\\&+ 23 \gamma_{268}+ \gamma_{311}+ \gamma_{349}+ \gamma_{359}- \gamma_{360}+ \gamma_{361}+ \gamma_{363}\\&- \gamma_{367}+ \gamma_{394}+ \gamma_{406}+ \gamma_{407}- \gamma_{64}+ \gamma_{67}- \gamma_{73}- 4 \gamma_{77}\end{aligned}\right)\\&+ \gamma_{241} \left(\begin{aligned}& 57 \varsigma  \gamma_{12} \gamma_{31}+ 82 \varsigma  \gamma_{13} \gamma_{35}+ 47 \varsigma  \gamma_{21} \gamma_{31}- 17 \gamma_{0} \gamma_{108}\\&+ 18 \gamma_{0} \gamma_{35}+ 14 \gamma_{100} \gamma_{12} \gamma_{13}+ 3 \gamma_{100} \gamma_{12} \gamma_{31}\\&+ 51 \gamma_{100} \gamma_{13} \gamma_{21}+ 7 \gamma_{100} \gamma_{13} \gamma_{33}- 17 \gamma_{100} \gamma_{23}\\&+ 14 \gamma_{12} \gamma_{31}- \gamma_{133}- \gamma_{147}+ \gamma_{152}- 5 \gamma_{193}+ 11 \gamma_{21} \gamma_{31}\\&- \gamma_{210}+ \gamma_{224}- 4 \gamma_{229}- 5 \gamma_{23}- 3 \gamma_{230}- 33 \gamma_{232}- 4 \gamma_{235}\\&- 11 \gamma_{236}- 27 \gamma_{237}- \gamma_{240}+ \gamma_{26}- \gamma_{273}- \gamma_{301}+ \gamma_{319}\\&- \gamma_{320}+ \gamma_{34}- \gamma_{352}- \gamma_{354}- 15 \gamma_{36}- \gamma_{371}- \gamma_{372}- \gamma_{373}\\&- \gamma_{374}- \gamma_{375}+ \gamma_{376}+ \gamma_{377}- \gamma_{379}- \gamma_{382}- \gamma_{39}+ \gamma_{399}\\&+ \gamma_{402}+ \gamma_{403}- \gamma_{405}- \gamma_{408}- \gamma_{409}- \gamma_{41}- 3 \gamma_{49}- \gamma_{52}\\&- 96 \gamma_{66}- 142 \gamma_{69}- 9 \gamma_{72}+ \gamma_{83}- 8 \gamma_{88}\end{aligned}\right)\\&- \gamma_{389} \gamma_{93}\end{aligned}\]

\[\mathsf A_{4,4}=\begin{aligned}&- \eta  \gamma_{187}+ \gamma_{101} \gamma_{255}+ 9 \gamma_{107}+ \gamma_{128}- 24 \gamma_{138}- 21 \gamma_{164}\\&- \gamma_{166}- 9 \gamma_{179}+ \gamma_{206}- 29 \gamma_{218}+ \gamma_{24}+ \gamma_{246} \gamma_{388}\\&- 42 \gamma_{269}+ \gamma_{295}+ \gamma_{298}+ \gamma_{326}+ \gamma_{369}+ \gamma_{381}- \gamma_{383}\\&- \gamma_{384}+ \gamma_{390}- \gamma_{401}+ \gamma_{404}+ \gamma_{409}+ 28 \gamma_{49}+ 44 \gamma_{66}\\&+ 30 \gamma_{69}+ 28 \gamma_{71}+ 26 \gamma_{73}\end{aligned}\]

\subsubsection*{Row 5}
\[\mathsf A_{5,1}=- \gamma_{140} \gamma_{411} \gamma_{92} \left(\varsigma  + \gamma_{141}\right) \left(\zeta  + \gamma_{366} + \gamma_{410} + \gamma_{77}\right)\]

\[\mathsf A_{5,2}=\begin{aligned}&- \gamma_{140} \gamma_{183} \gamma_{5} \left(\begin{aligned}& 3 \varsigma  \gamma_{12} \gamma_{13}+ 5 \varsigma  \gamma_{13} \gamma_{21}+ 2 \zeta  \gamma_{21} \gamma_{412}- \zeta  \gamma_{281}\\&+ 2 \zeta  \gamma_{35} \gamma_{412}- \gamma_{100} \gamma_{291}+ 2 \gamma_{12} \gamma_{13}- \gamma_{184}- \gamma_{247}\\&- \gamma_{285}- \gamma_{310}- \gamma_{338}+ \gamma_{35} \gamma_{412}- \gamma_{355}- \gamma_{395}- \gamma_{408}\\&- \gamma_{414}- \gamma_{71}\end{aligned}\right)\\&- \gamma_{142} \left(\begin{aligned}&- \gamma_{106}- \gamma_{107}+ 2 \gamma_{108}+ 5 \gamma_{111}+ \gamma_{116}- 3 \gamma_{117}\\&- \gamma_{15} \gamma_{412}+ \gamma_{171}+ \gamma_{178}+ \gamma_{189}+ \gamma_{328}+ \gamma_{364}+ \gamma_{365}\\&+ \gamma_{397}+ \gamma_{407}+ \gamma_{412} \gamma_{47}+ \gamma_{414}+ \gamma_{415}- \gamma_{416}+ \gamma_{418}\\&+ \gamma_{419}+ \gamma_{43} \gamma_{50}- \gamma_{75}\end{aligned}\right)\end{aligned}\]

\[\mathsf A_{5,3}=\begin{aligned}&- \gamma_{204} \left(\begin{aligned}&- \eta  \gamma_{412}- \zeta  \gamma_{100}- \gamma_{10} \gamma_{412}+ \gamma_{196}- 3 \gamma_{228}+ \gamma_{287}+ \gamma_{3}\\&+ \gamma_{345}+ \gamma_{368}+ \gamma_{392}+ \gamma_{400}+ \gamma_{406}+ \gamma_{410}+ \gamma_{413}+ \gamma_{415}\\&+ \gamma_{417}+ \gamma_{420}+ \gamma_{422}+ \gamma_{423}\end{aligned}\right)\\&- \gamma_{241} \gamma_{411} \left(\begin{aligned}&- \gamma_{228}- 3 \gamma_{268}+ \gamma_{307}+ 3 \gamma_{362}- 3 \gamma_{371}+ \gamma_{412} \gamma_{416}\\&+ \gamma_{418}+ \gamma_{421}+ \gamma_{422}- 3 \gamma_{74}\end{aligned}\right)\end{aligned}\]

\[\mathsf A_{5,4}=\begin{aligned}& 5 \varsigma  \gamma_{50}+ \gamma_{0} \gamma_{100} \gamma_{263}+ \gamma_{0}+ \gamma_{100} \gamma_{129}- 5 \gamma_{105}\\&- \gamma_{113} \gamma_{417}+ \gamma_{119}+ \gamma_{120}+ \gamma_{143}- \gamma_{157}- \gamma_{165}+ \gamma_{186}\\&+ \gamma_{261}+ \gamma_{264} \gamma_{29}+ \gamma_{266} \gamma_{412}- \gamma_{272} \gamma_{412}- \gamma_{315}\\&+ \gamma_{338}+ \gamma_{351}+ \gamma_{396}- \gamma_{412} \gamma_{424}+ 2 \gamma_{417}+ \gamma_{419}+ \gamma_{421}\\&+ \gamma_{423}+ \gamma_{424}- \gamma_{59}\end{aligned}\]
\endgroup

\clearpage
\section{The classical deformation gradients and pressure}\label{supp:classical}

The parameters and coordinates are those of
Section~\ref{sec:classicalmodels}. The gradient matrices below
use orthonormal cylindrical bases when a cylindrical coordinate is
present. The first row is radial and the third row axial.
For the homogeneous family $F=L\in\mathrm{SL}(3,\mathbb R)$.
For the block and straightening families, respectively,
\[
 F_1=\begin{pmatrix}\alpha/r&0&0\\0&\beta r&0\\0&\gamma&1/(\alpha\beta)\end{pmatrix},
 \qquad
 F_2=\begin{pmatrix}\alpha R&0&0\\0&\beta/R&0\\0&\gamma/R&1/(\alpha\beta)\end{pmatrix}.
\]
Their determinants are one. For the cylindrical family,
\[
 F_3=\begin{pmatrix}
 AR/r&0&0\\0&kr/R&\delta r\\0&\kappa/R&\gamma
 \end{pmatrix},\qquad
 \det F_3=A(k\gamma-\delta\kappa)=1.
\]
For the spherical family, after the fixed orthogonal map $Q$,
the gradient has radial value $\varepsilon R^2/r^2$ and tangential
values $r/R,r/R$. Its signed determinant is
$\varepsilon\det Q=1$. For the logarithmic cylindrical family,
\[
 F_5=\begin{pmatrix}\alpha&0&0\\\alpha\beta&\alpha k&0\\0&0&\gamma\end{pmatrix},
 \qquad\det F_5=\alpha^2k\gamma=1.
\]
Thus each matrix has the prescribed signed Jacobian.

For the block and cylindrical families, $B=FF^{\top}$ and
$c=B^{-1}$ have no radial off-diagonal entries and their physical
components depend only on $r$. The local pressure derivative is
\begin{equation}\label{supp:radial-pressure}
 \frac{dp_W}{dr}
 =2\frac{d}{dr}(W_1B_{rr}-W_2c_{rr})
  +\frac2r\{W_1(B_{rr}-B_{\theta\theta})
                       -W_2(c_{rr}-c_{\theta\theta})\}.
\end{equation}
For the spherical family the final factor $2/r$ becomes $4/r$.
For straightening, the tensors have no first-coordinate
off-diagonal entries and depend only on $x$; the pressure is
$p_W=2(W_1B_{xx}-W_2c_{xx})$ up to a constant.

The logarithmic cylindrical family has constant principal invariants.
In the current cylindrical frame,
\[
 B=\begin{pmatrix}\alpha^2&\alpha^2\beta&0\\
       \alpha^2\beta&\alpha^2(\beta^2+k^2)&0\\0&0&\gamma^2\end{pmatrix},
 \quad
 c=\begin{pmatrix}
 (\beta^2+k^2)/(\alpha^2k^2)&-\beta/(\alpha^2k^2)&0\\
 -\beta/(\alpha^2k^2)&1/(\alpha^2k^2)&0\\0&0&\gamma^{-2}
 \end{pmatrix}.
\]
Its local pressure one-form is
\[
 dp_W=2\{W_1(B_{rr}-B_{\theta\theta})-W_2(c_{rr}-c_{\theta\theta})\}\,d\log r
           +4(W_1B_{r\theta}-W_2c_{r\theta})\,d\theta.
\]
The coefficients are constant for each fixed energy $W$.
These formulas give the pressure outputs used in
Proposition~\ref{completion:sufficiency}.

\restoregeometry
\setcounter{section}{\suppSavedSection}
\setcounter{equation}{\suppSavedEquation}
\setcounter{table}{\suppSavedTable}
\endgroup
